% 本文件由 LaTeX 排版 Agent 根据有效 translation.json 自动生成。
% author=Tertullian; work_id=0306; title=致外邦人

\chapter{致外邦人}

% structure: p0001 source=h1 target=omitted reason=page-title-already-rendered-by-parent

卷一\newline{}卷二\par

\SourceRecord{Translated by Peter Holmes.；Ante-Nicene Fathers；Vol. 3.；Edited by Alexander Roberts, James Donaldson, and A. Cleveland Coxe.；Buffalo, NY: Christian Literature Publishing Co.,；1885.；<http://www.newadvent.org/fathers/0306.htm>.}

% structure: p0001 source=h1 target=section reason=page-title

\section{\WorkTitle{致外邦人（卷壹）}}

% structure: p0002 source=h2 target=subsection reason=relative-heading-rank; base=h2

\subsection{第一章——外邦人对基督徒的仇恨是不义的，因为这基于可责的无知。}

你们的无知，一方面谴责了你们的不义，一方面又为它开脱，一个明证立刻可见：所有曾与你们一样无知并仇恨（基督宗教）的人，一旦认识了它，便不再仇恨，因为他们不再无知；不仅如此，他们自己反而变成了从前所仇恨的，并开始仇恨他们从前所是的。的确，你们日复一日地抱怨基督徒人数增加。你们常喊说，\Quote{国家（被我们）围困；基督徒在你们的田里、在你们的营中、在你们的岛屿上}。你们把它当作灾祸而悲伤，说 \Quote{每个性别、每个年龄——简言之，每个阶层——都从你们那里转向我们}；然而即使在此之后，你们也不曾用心反省，这里面是否隐藏着某种好处。你们不容许自己有那些可能太过真实的猜疑，也不喜欢那些可能过于接近目标的尝试。这是唯一一种使人的好奇心变得迟钝的情况。你们乐于对别人欣喜发现的事情一无所知；你们宁愿不知道，因为你们如今怀有仇恨，仿佛知道（一旦认识，）你们的仇恨就必然终止。可是，如果仇恨没有正当理由，那么停止过去的不义必将被证明是最好的做法。然而，如果确有原因存在，仇恨就不会减少，反而会因意识到其正义而更加累积；除非，想必是，你们羞于改正过失，或不愿摆脱罪责。我很清楚你们通常用什么话来回应我们迅速增长这一论点。你们说，\Quote{那使许多人改宗并赢得他们支持的事物，不可仓促地视为好事}。我知道人心多么容易走上邪恶之路。多少人抛弃了正直的生活！多少人逃往相反的方向！确实很多，甚至非常多，正如末日临近。但这样的比较在运用上并不公平；因为众人都一致如此看待恶人，以至于连有罪者自己——他们站错了边，离开对善的追求而转向邪道——也不敢把恶辩护为善。卑鄙的事令他们恐惧，不敬的事令他们羞耻。总之，他们急于隐藏，回避公开，被抓住时战栗；被控告时否认；即使受刑，也不轻易或总是承认（罪行）；至少，被定罪时他们悲伤。他们为过去的生活责备自己；他们甚至把从无辜到邪恶品性的改变归因于命运。他们不能说那不是一件坏事，因此不肯承认那是自己的行为。然而，基督徒的情况与这有何相似？无人羞耻，无人懊悔，除了为从前的（罪）。若因（宗教）被人指点，他反以为荣；若被拖去受审，他不反抗；若被控告，他不辩护。被问及时，他承认；被定罪时，他喜乐。这是怎样的一种恶，在其中恶的本性竟停止了？\par

% structure: p0004 source=h2 target=subsection reason=relative-heading-rank; base=h2

\subsection{第二章——外邦人在审判基督徒时歪曲了判决。他们若完全取消审判程序会更一致。戴尔都良带着极大的愤慨力陈这一点。}

在这种情况下，你们实际上采用了与通常审判罪犯的司法程序相反的方式；因为当罪犯被提审时，如果他们否认指控，你们就用酷刑逼供。然而，当基督徒毫不勉强地承认时，你们却施加酷刑，诱使他们否认。这是何等大的乖谬，你们抗拒招供，\Emph{并}改变酷刑的用途，迫使坦承指控的人躲避，而迫使不愿承认的人否认？你们主审本为逼出真相，却唯独向我们要求谎言，好让我们宣称自己不是我们所是的。我想你们并不希望我们是坏人，因此热切地想把我们从那种品格中排除出去。诚然，你们将别人置于拉肢架和绞刑架上，好让他们否认他们所负的恶名。当他们否认（对他们的指控）时，你们不相信他们；但对我们，一旦否认，你们立即相信。如果你们确信我们是最有害的人，那么，为什么在对我们的诉讼中，你们对待我们的方式与对待其他罪犯不同呢？我不是说你们无视起诉或否认（因为你们的惯例并不是在没有起诉和辩护的情况下仓促定罪）；但以谋杀犯的审判为例，当一个人承认自己是杀人犯时，案件并不会就此了结，调查也不会满足。无论他的供认多么彻底，你们并不轻易相信他；不仅如此，你们还调查附带情况——他犯过多少次谋杀；用什么武器，在什么地方，抢了多少财物，同谋，\Emph{以及}事后从犯（罪行是如何实施的）——为的是不遗漏关于罪犯的任何细节，并且备好一切手段以达成真正的判决。相反，在我们身上，你们相信我们犯有更凶残、更多的罪行，却用更简短、更轻微的措辞来拟定起诉。我想你们并不在意用指控来加重那些你们急于摆脱的人的负担，或者你们认为没有必要调查那些你们已经知道的事情。然而，你们强迫我们否认那些你们有最明确证据的指控，这就更加乖谬了。但是，确实，如果你们出于对我们的仇恨，干脆免去一切司法程序，竭尽全力不催逼我们说\Quote{不，}，从而不得不宣告你们所仇恨的对象无罪；反而要我们承认所有各项加给我们的罪行，这样你们的怨恨就可以因我们惩罚的累积而更好地满足，当人们知道我们每个人可能参加过多少次那些宴席，在夜幕掩护下犯过多少次乱伦时，岂不更加一致！我在说什么？既然你们根除我们团体的调查必须广泛进行，你们应当把调查扩大到对我们的朋友和同伴。让我们的杀婴者和（我们可怕宴席的）厨师被带出来——是的，还有服务于我们（乱伦）婚礼的狗；那么（我们审判的）事务就没有缺点了。即使对挤满观赏的人群也会增添趣味；因为当一个人必须与一个吞食了一百个婴儿的对手在竞技场上战斗时，他们将会怀着多大的热切涌入剧场啊！因为既然这些可怕而骇人的罪行被报道与我们有关，它们当然应当被揭露，以免它们显得难以置信，并且公众对我们的憎恶开始冷却。因为大多数人不轻易相信这种事情，一想到我们的本性竟然渴望野兽的食物，就感到极其厌恶，而自然已经阻止了这些野兽与人类的一切交合。\par

% structure: p0006 source=h2 target=subsection reason=relative-heading-rank; base=h2

\subsection{第三章——基督徒的大罪在于他们的名字本身。名字得到辩护。}

既然你们在其他案件中追查远为轻微的指控时，最为审慎与执着，却在像我们这样可怕且罪大恶极的案件中——即使用\Emph{不虔敬}一词也嫌太轻——放弃你们的职责，拒绝听取宣认（这本应是主理司法程序者始终重视的过程），且未能进行充分调查（这本应是起诉定罪者所应为之事），那么显而易见，归咎于我们的罪行并非在于任何有罪行为，而完全在于我们的\Emph{名称}。实际上，若任何真实罪行能明确地加诸于我们，其名称本身便会定我们的罪——若适用的话，以至于会如此对我们分别宣判：让那个凶手，或那个乱伦犯，或我们被指控的任何罪行，被处决、被钉十字架或被投给野兽。然而，你们的判决仅意味着一个人承认自己是基督徒。没有罪行的名称指向我们，只有名称的罪行。而这实际上恰恰就是对我们怀有的全部憎恨。名为缘由：某种因你们的无知而强化的神秘力量攻击它，以至于你们不愿确切地知道那些你们确定一无所知的事；因此，你们也不相信未经证实的事，并且为了避免它们轻易被驳倒，你们拒绝进行调查，致使这个可憎的名称在（真实）罪行的推定下受到惩罚。因此，为了使案件脱离这个冒犯的名称，我们被迫否认它；然后在我们否认后，我们被宣告无罪，过去的一切完全赦免：我们不再是凶手，不再是乱伦者，因为我们失去了那个名称。但由于这点已在别处讨论，请你们明确告诉我们，为何你们要追捕这个名称直至灭绝？一个名称有什么罪行、什么冒犯、什么过失？因为规则禁止你们提出（任何人的）罪行，这些罪行没有法律诉讼提出，没有起诉书详细说明，没有判决书列举。在任何提交给法官、针对被告进行调查、由他回应或否认、并从法官席上引用的案件中，我承认一项法律指控。因此，关于一个名称的价值，无论名称可能被指控何种冒犯，无论词语可能应受何种弹劾，我个人认为，连一句抱怨也不应归于一个词语或名称，除非它听起来野蛮，或带有厄运，或是不庄重的，或是对说话者不体面，或是对听者不愉快。这些在（单纯的）词语和名称中的罪行，就像野蛮的词语和短语一样，有其错误、语病和可笑的比喻。然而，\OriginalTerm{基督徒}{Christian}这个名称，就其含义而言，带有傅油的意思。即使你们通过错误发音称我们为\Quote{\OriginalTerm{克勒斯提安}{Chrestians}}（因为你们连这个著名名称的发音都不确定），你们事实上也在含糊不清地说出愉悦和良善的含义。因此，你们是在诽谤无辜之人，甚至是我们所怀的无害名称，这个名称既不碍口，也不刺耳，也不伤害任何人，也不粗鲁地对待我们的国家，它是一个美好的希腊词，正如许多其他词一样，音义皆悦耳。确实，确实，名称不是应受刀剑、十字架或野兽惩罚的东西。\par

% structure: p0008 source=h2 target=subsection reason=relative-heading-rank; base=h2

\subsection{第四章——在基督徒身上受憎恨的真理；古时在\Person{苏格拉底}身上也是如此。基督徒的美德。}

但你们说，这教派以其创始人之名而受惩罚。首先，诚然，一个教派以其创始人之名来区分是公平且常见的做法，因为哲学家以其师之名被称为\Person{毕达哥拉斯}派和\Person{柏拉图}派；同样，医生被称为\Person{埃拉西斯特拉图}派，文法家被称为\Person{阿里斯塔尔科}派。因此，如果一个教派因其创始人恶劣而名声不佳，它便作为恶名的传统承担者而受惩罚。但这是一种轻率的假设。首先应当了解创始人是什么样的人，以便理解其教派，而不是从教派出发阻碍对创始人品格的探究。但在我们的情形中，由于你们必然对教派无知，因你们不知其创始人，或者由于你们不审视其教派而未能公正地考察创始人，你们只抓住名字不放，仿佛以此既诋毁教派又诋毁创始人，而对两者一无所知。然而，你们却公开允许你们的哲学家有加入任何学派并以创始人名字作为自己称号的权利，而无人激起对他们的仇恨，尽管他们公开和私下里对你们的习俗、礼仪、仪式和生活方式吠出最尖刻的言辞，如此藐视法律、不敬人物，甚至对皇帝本人也放肆妄言而免受惩罚。然而，正是真理——它令世界如此烦恼——为这些哲学家所追求，但为基督徒所拥有；因此，拥有真理的人更令人不悦，因为追求真理者玩弄它，拥有真理者持守它。例如，\Person{苏格拉底}因为在其探索中最接近真理的那一面——即摧毁你们的神明——而被定罪。虽然当时世上尚无基督徒之名，但真理始终在受谴责。现在你们不会否认他是个智者，因为你们自己的皮提亚（神）曾为他作证。他说：\Quote{\Person{苏格拉底}是众人中最有智慧的。}真理战胜了\Person{阿波罗}，迫使他甚至做出对自己不利的宣告，因为他承认自己不是神，同时断言那否定诸神的人是最有智慧的。然而，按你们的原则，他因否定诸神而更少智慧，尽管事实上他因这否定而更有智慧。同样地，你们习惯于这样谈论我们：\Quote{\Person{路奇乌斯·提提乌斯}是个好人，只是他是基督徒；}另一个人说：\Quote{我真奇怪，像\Person{盖尤斯·塞尤斯}这样可敬的人竟成了基督徒。}照他们愚蠢的盲目，人们赞美他们所知道的，指责他们所不知的；他们以所不知的来败坏所知道的。没有人想到要思考：一个人是否因为是基督徒而不好、不智慧，或者因为他智慧、好而成为基督徒，尽管人类行为中更常见的是以明显之事确定晦暗之事，而不是以晦暗之事来预判明显之事。有些人惊讶于那些在拥有这个名字之前他们已知为反复无常、卑鄙\Emph{或}邪恶的人突然转变向善；然而他们更善于对此惊讶而非达致此转变；另一些人则固执地争斗，与自己的最佳利益为敌，而这些利益本可通过与那可憎的名字交往而获得。我知道不止一个丈夫，先前为妻子的品行焦虑，连老鼠爬进卧室都忍不住怀疑呻吟，一旦发现她们新近勤勉、前所未有地专注于家务的原因后，竟将妻子完全借给他人，放弃一切嫉妒，宁愿做母狼的丈夫也不做基督徒女人的丈夫：他们能放任自己于违背自然的滥用，却不能容忍妻子改过向善！一位父亲剥夺了儿子的继承权，而之前他已不再指责儿子。一位主人将他的奴隶送进惩戒所，而之前他发现这奴隶对他不可或缺。一旦发现他们是基督徒，他们就希望他们再次成为罪犯；因为我们的纪律本身就带着自身的证据，除了我们的良善外，我们不被任何事所出卖，正如恶人也因自身邪恶而显明。否则，为何唯独我们，违背自然之理，因我们的善而被烙上极恶的印记？因为我们显出了什么标记，除了首要的智慧——它教导我们不崇拜人手所造的轻浮作品；节制——使我们不贪图他人之物；贞洁——我们甚至不以眼神玷污它；怜悯——促使我们帮助穷人；真理本身——它使我们冒犯他人；以及自由——我们甚至为之学习死亡？凡愿了解基督徒是谁的人，必须用这些标记来发现他们。\par

% structure: p0010 source=h2 target=subsection reason=relative-heading-rank; base=h2

\subsection{第五章——任何假基督徒的言行不一并不比一片过云遮蔽夏日天空更能定真正基督门徒的罪。}

至于你们称我们为最可耻的一群，完全沉溺于奢侈、贪婪和败坏，我们不否认这对某些人是真的。然而，这足以证明我们的名称，即不能对所有的人这样说，甚至不能对我们中的大多数这样说。即使在最健康、最纯洁的身体上，也难免会长疮、生疣或斑点瑕疵。连天空本身也不是完全晴朗，没有一丝薄云点缀。脸上的一点瑕疵，因为它在如此显眼的部分显而易见，反而证明了整个肤色的纯净。大多数人的善良由这微小的瑕疵得到了充分证明。但是，虽然你证明我们中间有些人是邪恶的，你并不能因此证明他们是基督徒。请查考并看看，是否有哪个派别会把局部的缺点归咎为整体的污点。你们自己在谈论中习惯于贬低我们说：\Quote{为什么某某人奸诈，而基督徒却如此克己？为什么残忍，而基督徒却如此慈悲？}你们这样作证，说明这不是基督徒的品格，因为你们以反驳的方式问，那些被称为基督徒的人怎么会有这样的性情。指责和名称之间、意见和真理之间有很大区别。因为名称被指定，正是为了在单纯的称谓与实际状况之间划定适当界限。实际上，有多少被称为哲学家的人，却并不履行哲学的信条？所有人都因其所宣称的职业而拥有此名；但他们只拥有称号，却没有职业的卓越，他们以名称的浅薄伪装玷污了真实事物。人并不因为被说成具有某种品格就立刻具有那种品格；但若他们没有，这样称呼他们是徒劳的：他们只欺骗那些将现实附着于名称的人，而事实上是名称与事实的一致性决定了名称所隐含的状况。然而，这种可疑类型的人不同我们聚集，也不属于我们的共融；因他们的过失，他们再次成为你们的人，因为我们甚至不愿与那些被你们的暴力和残酷逼迫背弃的人交往。当然，我们更愿意将那些不情愿放弃我们纪律的人包括在我们中间，而不是那些故意背教的人。然而，你们无权称他们为基督徒，因为基督徒自己否认他们这个名称，而他们也未曾学会否认自己。\par

% structure: p0012 source=h2 target=subsection reason=relative-heading-rank; base=h2

\subsection{第六章——基督徒的无辜并未因反对他们的不义法律而受损害。}

每当我们的这些陈述和回答——真理自发地提示——压迫并约束着你们的良心，而这良心正是它自身无知的见证，你们就匆匆忙忙地逃向那可怜的避难祭台，即法律的权威，因为，如果立法者没有充分考虑过这冒犯性派别的应得之罪，这些法律当然不会惩罚它。那么，是什么阻止了执行法律的人进行类似的考虑呢？因为对于其他所有同样被法律禁止和惩罚的罪行，除非经过正规程序追诉，否则不会施加刑罚。例如，以杀人犯或奸淫者为例。即便罪行的性质对所有人都是明显的，也仍会就犯罪细节下令进行调查。基督徒所行的任何错事都应当被揭露出来。没有法律禁止进行调查；相反，调查是为了法律的利益。因为，如果你因为未能察觉你所需要遵守的是什么而抵消了防范的谨慎，你又如何通过防范法律所禁止的事情来遵守法律呢？任何法律都不应将自身正义的知识保留给自己，而应将其归于那些要求服从的人。然而，当法律拒绝证明自身时，它就变成可疑的对象。因此，很自然地，只要人们对其目的和含义仍然无知，反对基督徒的法律就被认为是公正的，值得尊重和遵守；但当这一点被察觉时，它们的极大不义就被发现，它们连同其酷刑工具——刀剑、十字架和狮子——就被理所当然地以憎恶加以拒绝。不公正的法律得不到尊重。然而，在我看来，你们中间有一种怀疑，认为这些法律中有一些是不公正的，因为你们没有一天不在通过新的审议和决定来修改它们的严厉和不义。\par

% structure: p0014 source=h2 target=subsection reason=relative-heading-rank; base=h2

\subsection{第七章——基督徒被诽谤。对名声的讽刺性描述；其欺骗和对基督徒的恶毒诽谤的详细描述。}

你们会对我们说：怎么会发生这样的事，以致这种性格被加于你们，或许还因为这种指责而证明了立法者的正当？让我反过来问：他们当时有什么凭据，或者你们现在有什么凭据，证明这种指责的真实性？（你们回答：）名声。那么，这难道不是——\par

\Quote{名声是恶，没有什么比它更快？}\par

那么，如果流言永远都是真的，为什么它被认为是\Emph{一种灾难}？它从未停止说谎；即使在它报告真相的那一刻，它也并非没有说谎的意愿，以至于不会通过添加、减少或混淆各种事实的方式将虚假与真实交织在一起。事实上，它的状况就是如此：它只有通过说谎才能继续存在。因为它的生命仅仅在于它无法证明任何事情。一旦它证明自己是真实的，它就垮台了；而且，就好像它的新闻报道职责已经结束，它便离开岗位：此后，这件事就被视为事实，并以那个名称存在。因此，没有人会说，举例来说，\Quote{这件事发生在\Place{罗马}的报道}，或\Quote{他获得了一个省份的传闻}；而是说\Quote{他获得了一个省份}以及\Quote{这件事发生在\Place{罗马}}。除非在不确定的情况下，没有人会提到传闻，因为没有人能确定一个传闻，只能确定确切的知识；只有傻瓜才会相信传闻，因为没有智者会相信不确定之事。无论一个消息传播得多么广泛，它必定最初是从某一个人的口中开始的；然后它不知怎地传到那些传递它的耳朵和舌头上，从而掩盖了最初的无辜错误，以至于没有人考虑最初传播它的那个嘴巴是否散布了谎言——这种情况常常发生在竞争心态、怀疑倾向、甚至编造新闻的乐趣中。然而，好在时间揭示一切，正如你们自己的谚语和格言所证明的；是的，就像自然本身所证明的那样，它如此安排，以至于没有什么是隐藏的，即使是名声没有报道的东西也是如此。看，现在你们对我们使用了什么样的证人：它至今无法证明它引发的传闻，尽管它有如此长的时间来让我们接受。我们的这个名称起源于\Person{奥古斯都}统治时期；在\Person{提庇留}时期，它被清晰而公开地教导；在\Person{尼禄}时期，它被无情地定罪，你们甚至可以从迫害者的性格来权衡它的价值和特性。如果那位君主是虔诚的人，那么基督徒就是不虔诚的；如果他是公义的，如果他是纯洁的，那么基督徒就是不公义和不纯洁的；如果他不是公敌，那么我们就是国家的敌人：我们是什么样的人，我们的迫害者自己就表明了，因为他当然惩罚那些对他产生敌意的事物。现在，尽管\Person{尼禄}统治下的所有其他制度都已被摧毁，但我们的这个制度却稳固地存留——这似乎是公义的，因为它与迫害的发起者不同。那么，从我们的生命开始到现在，还不到二百五十年。在此期间，有如此多的罪犯；如此多的十字架获得了不朽；如此多的婴儿被杀害；如此多蘸血的饼；如此多的蜡烛熄灭；如此多淫乱的婚姻。直到现在，仅仅是传闻在与基督徒作战。毫无疑问，它在人心的邪恶中找到了强有力的支持，并且在残酷野蛮的人中间更成功地散布谎言。因为你们越是倾向于恶意，就越容易相信邪恶；简而言之，人们更容易相信虚假的邪恶，而不是真实的良善。现在，如果非正义在你们心中为审慎调查传闻的真实性留下了任何空间，那么正义当然要求你们审查是谁可能在民众中传播了这个传闻，并从而流传到整个世界。因为我想，它不可能由基督徒自己传播，因为根据所有奥秘的宪法和法律，沉默的义务被强加。对于像加给我们的那些（奥秘）来说，更是如此，如果泄露，必然立即招致人们迅速报复的惩罚！因此，既然基督徒不是自己的告密者，那么必然是外人。现在我问，外人怎么可能了解我们，当即使是真正的和合法的奥秘也排除任何外人观看时，除非非法的奥秘不那么排他？那么，外人更可能既不知道（案情的真实情况），又编造（虚假的叙述）。也许我们的家仆偷听、透过缝隙和孔洞窥视，并偷偷获取了我们行事方式的信息。那么，当我们的仆人向你们告密时，我们该说什么？更好的（当然）是我们不被任何人出卖；但尽管如此，如果我们的行为如此残暴，当义愤甚至冲破家庭忠诚的所有纽带时，那不是更恰当吗？它怎么可能忍受那使心灵恐惧、使眼睛惊恐的事情？这也是一件奇妙的事：那个如此被不耐烦的激动所压倒以至于成为告密者的人，他难道没有同样想要证明（他所报告的事情）吗？那个听到告密者故事的人也不关心自己去看，因为毫无疑问，告密者证明他所报告的，和听者使自己相信他所听到的可信，两者的奖赏是相等的。但你们说（这正是所发生的事）：先是传闻，然后是证据的展示；先是道听途说，然后是亲眼目睹；此后，名声接受了它的使命。现在我必须说，这超越了所有惊奇：那种事情被一次性地发现和泄露，而它却不断重复发生，除非我们实际上已经停止了这些事情（如加给我们的）的重复。但我们仍然被称为同样的（冒犯性）名字，我们被认为仍然从事同样的行为，并且我们一天天增多；我们越多，就越成为憎恨的对象。憎恨随着其材料的增加而增加。现在，既然犯罪者的数量不断增长，告密者的数量为什么不保持同步？据我所知，甚至我们的生活方式也变得更加为人所知；你们知道我们集会的具体日子；因此，我们被围困、被攻击，甚至在我们的秘密集会中被囚禁。然而，有谁曾在我们中间发现过半腐烂的尸体？有谁在我们蘸血的饼上发现过咬痕？有谁曾通过突然的光亮侵入我们的黑暗，发现任何不洁的痕迹，我不说乱伦的痕迹？如果我们通过贿赂使自己免于被这样拖到公众面前，我们怎么仍然受到压迫？我们确实有能力根本不被这样逮捕；因为谁要是出售或购买关于罪行的信息，如果罪行本身不存在呢？但我何必贬低地提及陌生间谍和告密者，当你们指控我们那些我们肯定不会大肆宣扬的罪名时——要么是你们一听到就指责，如果我们事先向你们展示了它们，要么是你们自己发现之后，如果当时它们对你们隐藏的话？因为毫无疑问，当任何人想要加入这些奥秘时，他们的习惯是先去见圣礼的主持或父亲。然后他（对申请者）会说：\Quote{你必须带来一个婴儿，作为我们仪式的保证，要被献祭，以及一些要掰开并蘸在他血中的饼；你还需要蜡烛、拴在一起的狗来打翻它们，以及一些肉来激起狗。此外，一位母亲或一位姐妹对你也是必需的。但是，如果你两者都没有，又该怎么办？我想那样你就不能成为一个真正的基督徒。}现在，让我问问你们，这样的事情被外人报告时，能够被传播（作为对我们的指控）吗？这些人不可能理解他们没有参与的事情。过程的第一步是用诡计进行的；我们的宴席和我们的婚姻是由无知的人编造和详细描述的，他们从未听说过基督徒的奥秘。尽管他们后来不可避免地获得了一些关于它们的知识，但即使那时，也是作为由他们带入场景的他人来管理的事情。此外，外邦人知道司铎不知道的奥秘，这是多么荒谬！那么，他们自己保留所有这些，并视之为理所当然；因此，这些悲剧（比\Person{提厄斯忒斯}或\Person{俄狄浦斯}的更坏）完全没有暴露在光下，也没有传入公众耳中。即使是更贪婪的咬伤也不会从那些被接纳者（无论是仆人还是主人）的信誉中带走任何东西。然而，如果这些指控中没有一件能被证明是真的，那么那个宗教的伟大（有多么不可估量），它显然没有被这些巨大暴行的重担所压倒！哦，外邦人，你们值得同情，看哪，我们把我们的神圣体系所提供的应许摆在你们面前。它保证永生给那些跟随并遵守它的人；另一方面，它以永不熄灭之火所罚的永久惩罚来威胁那些不敬和敌对的人；而对两类人，都宣讲死人复活。我们现在不关心这些（真理）的教义，这些教义在适当的地方讨论。但同时，请相信它们，就像我们自己相信的那样，因为我想知道你们是否准备好像我们一样通过这样的罪行达到它们。来吧，无论你是谁，将你的剑刺入一个婴儿的体内；或者如果那是别人的职责，那么只消看着那有生命气息的造物在它尚未活过之前死去；至少，接住它新鲜的血来蘸你的饼；然后毫无节制地喂养自己；在这过程中，斜躺着。仔细区分你的母亲或你的姐妹可能铺床的地方；记好它们，以便当夜晚的阴影降临到它们身上时，当然要考验你们每个人的细心，你们不会错误地落到别人身上：如果你没有乱伦成功，你就要做补赎。当你完成这一切时，永生将为你存留。我想让你告诉我，你是否认为永生值得这样的代价。不，的确，你不相信；即使你相信，我断言你也不愿意付出（代价）；或者如果愿意，也会无法做到。但如果你们不能，为什么别人能？如果别人不能，为什么你们能？你们希望不罚和永恒要付出什么代价？你们认为这些（祝福）可以用任何代价买到吗？基督徒的牙齿与别人的不同吗？他们有更宽大的下颌吗？他们对乱伦的情欲有更强的神经吗？我想没有。我们只凭真理与你们在状况上不同，这就够了。\par

% structure: p0018 source=h2 target=subsection reason=relative-heading-rank; base=h2

\subsection{第八章。——借\Person{普萨美提克}的发现所说明的对基督徒的诽谤。对故事的驳斥。}

我们确实被称为\Quote{\OriginalTerm{第三种族}{third race}}的人。什么，狗脸种族？还是宽脚影？或地下的\OriginalTerm{对跖人}{Antipodes}？如果你们对这些名称有所理解，请告诉我们第一和第二种族是什么，以便我们了解这个\Quote{第三}。\Person{普萨美提克}以为他巧妙地发现了原始人。据说他将某些新生儿与所有人际交往隔离，并将他们托付给一位预先被割去舌头的乳母，以便他们完全远离一切人声，能在不听到说话的情况下形成自己的语言；这样，仅凭自身，就能指示出那第一个其语言由自然所规定的民族。他们的第一句话是\OriginalTerm{贝克科斯}{Bekkos}，这个词在\Place{弗里吉亚}语中的意思是\Quote{\Emph{面包}}：因此，弗里吉亚人被认为是最早的人类。但是，我们不妨做一个观察，以显示你们的信德如何更多地倾向于虚妄而非真理。那么，乳母在失去了如此重要的肢体——正是生命气息的器官，而且是从根部割断，喉部被割伤，即使外部受伤也能致命，脓血流回胸腔，长时间没有食物——之后还能存活，这真的可信吗？来吧，即使假设通过\Person{菲洛墨拉}的疗法她存活了下来，就像智者所设想的，他们认为哑巴不是因为割舌，而是因为羞愧脸红；如果在这种情况下她活着，她仍然能够发出一些模糊的声音。仅凭张口发出的尖锐不清的噪音，没有任何嘴唇的调节，即使没有舌头帮助，也能从喉咙中挤出。很可能，婴儿们很容易模仿这个声音，尤其是因为这是唯一的声音；只是他们做得稍微清晰一些，因为他们有舌头；然后他们赋予了它一个明确的意义。那么，就算弗里吉亚人是最早的种族，也不意味着基督徒是第三种族。因为在弗里吉亚人之后还有多少其他民族？不过要注意，你们称之为第三种族的人可能获得第一的地位，因为实际上没有一个民族不是基督徒。因此，无论哪个民族是最早的，现在仍然是基督徒。你们称我们为最晚的种族，然后又特别称我们为第三，这是荒谬的愚蠢。但是，我们被认为是第三，是就我们的宗教而言，而非我们的民族；序列是：罗马人、犹太人，以及随后的基督徒。那么，希腊人在哪里？或者，如果他们因迷信而被算作罗马人（因为罗马甚至从希腊借来了诸神），那么埃及人至少在哪里？因为据我所知，他们有自己独特的神秘宗教。现在，如果属于第三种族的人如此怪异，那么首先和其次在他们之前的人应该被认为是怎样的呢？\par

% structure: p0020 source=h2 target=subsection reason=relative-heading-rank; base=h2

\subsection{第九章。——基督徒并非公共灾难的原因：在基督教之前已有这样的灾祸。}

但我何以要惊讶于你们虚妄的指控呢？恶意与愚妄总是以同一天然形态合为一体、共同成长，并常在同一个错谬煽动者之下反对我们。诚然，我并不惊讶；因此，为论述所需，我将列举一些实例，好叫你们因列举出你们所陷入的愚妄而感到惊讶——当你们坚持认为我们是每一公共灾难或损害的起因时。若台伯河泛滥，若尼罗河仍在河床内，若天无风雨，或地有震动，若死亡肆虐，或饥荒降灾，你们立刻便喊道：\Quote{这是基督徒的过错！}仿佛敬畏真天主之人竟会害怕轻微之事，或至少害怕（除地震、饥荒或此类灾难之外的）任何事。我想，正是因为我们蔑视你们的神，才招致他们如此打击我们。正如我们已指出的，我们存在至今尚未满三百年；但在此之前，何等巨大的灾难降于全世界，降于其各个城市与省份！历史记载了何等可怕的战争，既有外战也有内战！何等瘟疫、饥荒、火灾、地裂与地震！那时，基督徒在哪里？当罗马国提供了如此多的灾难记录时，基督徒在哪里？当希拉岛、阿纳斐岛、提洛岛、罗得岛、凯阿岛因众多人口而荒芜时，基督徒又在哪里？或者，当柏拉图所说比亚细亚和非洲还大的那块陆地沉没于大西洋时？或者，当天火吞噬了沃尔西尼，来自自家山上的火焰吞灭了庞贝时？当科林斯海被地震吞没时？当全世界被洪水毁灭时？\Emph{那时}（我不说蔑视你们神的基督徒，而是）你们的神自身何在？他们被证明比那场大毁灭更晚出现，这由他们出生、旅居、埋葬及亲自建立的地点与城市便可知。否则他们不会存留至今，除非他们晚于那场灾难。如果你们不愿研读和反思这些历史见证——这些记录对你们和我们的影响不同——尤其是为使你们不必指责你们的神极端不义，因为他们竟因着蔑视他们的人而伤害崇拜他们的人——那么，当你们以他们为神，而他们却不区分你们与渎神之人的功过时，你们岂不证明自己也是错的吗？然而，若如有时虚妄所说，你们因在铲除我们的事上懈怠而招致你们神的惩罚，你们便已确定了他们软弱和微不足道；因为如果他们自己能做什么，就不会因你们在惩罚我们的事上拖延而向你们发怒——尽管你们实际上以另一种方式承认了同一件事：每当你们惩罚我们，仿佛是替他们报仇时。若一方利益由另一方维护，那么维护者大于被维护者。那么，神竟被凡人维护，这是何等羞耻！\par

% structure: p0022 source=h2 target=subsection reason=relative-heading-rank; base=h2

\subsection{第十章——基督徒并非唯一蔑视神的人。对神的蔑视常由异教官员表现。荷马使神成为可鄙者。}

现在，倾尽你们所有的毒液吧；用你们所有毁谤之箭射向我们这名号吧：我不再久留以驳斥它们；但当我们详述我们的全部教规时，它们不久便会钝化。现在，我确实满足于将这些箭从我们自己身上拔出，并掷回你们身上。你们以指控加于我们的创伤，我将证明也印在你们自己身上，好使你们被自己的剑和镖所击倒。现在，首先，当你们以背叛祖先制度的普遍罪名指控我们时，请反复思考，你们自己是否不也与我们共同承受这一指控。因为我观察你们的生活和习俗，看哪，我所发现的难道不是旧制度已被败坏，甚至被你们毁灭了吗？关于法律，我已说过，你们每天都在以新的法令和法规取代它们。至于你们生活方式的其余一切，你们从祖先那里做出了多么大的改变——在风格、衣着、仪仗、饮食，甚至言语上；因为你们驱逐古老样式，仿佛它冒犯了你们！无论在公共事务还是私人职责中，古制都被废除了；你们自己的权威取代了祖先的一切权威。诚然，你们永远赞美旧习俗；但这只给你们带来更大的耻辱，因为你们始终拒绝它们。你们的乖谬该有多大，竟认可祖先的制度（它们低劣得不能持久），而同时又拒绝你们所认可的对象！但即使是你们祖先的这份家传，你们似乎最忠实地守护和捍卫，并从中找到最强有力的理由来控告我们为违法者，且使你们对基督徒之名的仇恨获得全部生命——我指的是对诸神的敬拜——我将证明它正从你们自己那里遭受毁灭和轻蔑，丝毫不亚于我们——除非没有理由认为我们像你们一样是藐视诸神的人，因为没有人会藐视他所知完全不存在的东西。确实存在的东西可能被藐视。那\Emph{是}虚无的，不受任何损害。因此，对于那些视之为存在者的人来说，必然会产生影响它的损害。那么，指控就更加沉重地压在你们身上，你们相信有神，却同时藐视他们，既敬拜又拒绝，既尊敬又攻击。你也可以从这一考虑得出同样的结论：既然你们敬拜不同的神，各自敬拜一些，你们当然藐视那些你们不敬拜的神。偏爱一个不可能不经视另一个，选择不可能没有拒绝。从众多中选取某个\Emph{一}个者，已经轻视了他未选取的那一个。但如此众多且伟大的神不可能被所有人敬拜。那么，你们从一开始就必然行使了你们的藐视，因为你们当时并不害怕如此安排，以至于所有神都不能成为所有人敬拜的对象。因为你们那些非常明智和谨慎的祖先，你们不知道如何废除他们的制度，尤其是在敬神方面，他们自己被发现是不虔敬的。我大错特错，如果他们不曾有时下令，任何将军在战斗中立愿修建的神庙，在元老院批准之前不得奉献；正如\Person{马尔库斯·埃米利乌斯}的案例，他曾向神\OriginalTerm{阿尔布尔努斯}{Alburnus}许愿。那么，将神祇的荣耀置于人类判断的意志和喜好之下，以至于非经元老院允许就不能有神，这岂不是公然最大的不虔敬，乃至最大的侮辱？监察官们曾多次未经人民同意就毁掉一个神。父神\OriginalTerm{巴克斯}{Bacchus}及其全部礼仪，确实被执政官依据元老院权威，不仅逐出城市，而且逐出整个\Place{意大利}；而\Person{瓦罗}告诉我们，\OriginalTerm{塞拉皮斯}{Serapis}、\OriginalTerm{伊西斯}{Isis}、\OriginalTerm{阿波克拉特斯}{Arpocrates}和\OriginalTerm{阿努比斯}{Anubis}也被逐出\Place{卡皮托利山}，且元老院推倒的祭坛仅因民众暴力才得以恢复。然而，\Person{加比尼乌斯}执政官在随后的一月第一天，虽然因聚集的人群而迟迟同意某些祭祀，因为他未能就\OriginalTerm{塞拉皮斯}{Serapis}和\OriginalTerm{伊西斯}{Isis}作出决定，但他仍认为元老院的判决比群众的喧哗更有力，并禁止修建祭坛。那么，在你们自己的祖先中，你们就有了轻视诸神的基督徒，即使不是其名，也至少有其行径。然而，即使你们在敬礼神时并无背叛他们的罪责，我仍发现你们在迷信和不虔敬方面都保持了一贯的进步。因为你们被发现是多么更无信仰！你们有家神\OriginalTerm{家神}{Lares and Penates}，通过家庭祝圣而拥有：你们自己和家仆甚至亵渎地将他们踩在脚下，为了需要或任性而叫卖或典当他们。这种无礼的亵渎，如果不是针对你们更低级的神祇，或许还可原谅；但事实只是更无礼。然而，你们的私家家神在这些侮辱中得到一些安慰，即你们对待公共神祇时更加侮辱和无礼。首先，你们为他们做广告拍卖，提交公开出售，出给最高出价者，每五年将它们与你们的税收一起敲槌出售。为此，你们频繁出入\OriginalTerm{塞拉皮斯}{Serapis}神殿或\Place{卡皮托利山}，在那里举行销售，签订合同，仿佛它们是市场，伴随着众所周知的拍卖官声音和相同的财务官征收。如今，土地因赋税而变得便宜，人头税使人的价值降低（这些是众所周知的奴隶制标志）。但诸神，缴纳的贡赋越多，就越神圣；或者更确切地说，越神圣，缴纳的贡赋越多。他们的威严被转化为交易物品；人们以宗教为业；诸神的神圣因销售和合同而贫穷。你们将你们神殿的地面、祭坛的通道、供物、祭品都变成了商品。你们出卖诸神的全部神性。你们不允许他们被无偿敬拜。拍卖官比祭司更需要修缮。\par

若要检验你们轻视的程度，仅凭你们无耻地利用你们的神明牟利，还不够；你们不仅不给予他们尊荣，还用某种侮辱贬低你们给予他们的那一点点。其实，你们为尊崇你们的神明所做的一切，何尝不是同样施之于你们的亡者？你们为神明建造庙宇，也为亡者建造庙宇；你们为神明筑造祭台，也为亡者筑造祭台；你们在两者之上刻写同样的题铭；你们为他们的雕像勾勒同样的轮廓——按他们各自的才德、职分或年龄：你们将\Person{萨图尔努斯}塑成老人，\Person{阿波罗}塑成无胡须的青年；你们从\Person{狄安娜}塑造一位贞女；在\Person{玛尔斯}身上你们祝圣一名士兵，在\Person{武尔坎}身上祝圣一名铁匠。因此，你们为亡者献上同样的牺牲、焚烧同样的香料，如同为神明所做，就不足为奇了。那种将各类亡者与神明等量齐观的傲慢，能找到什么借口？甚至你们的君王也被分派了司祭的职务、神圣的仪式、战车、轿车，以及\OriginalTerm{坐席祭}{solisternia}和\OriginalTerm{躺席祭}{lectisternia}的尊荣，还有节日和竞技。诚然，天堂为他们敞开，但这仍然是对神明的侮辱：首先，因为即便他们出生后得以成为神明，将其他存在与他们并列也绝不可能得体；其次，因为目睹那人被提到天上的见证人，若不是因他和那些容许伪证的人对所发誓的对象心存轻视，就不会在民众面前如此随意而明显地发假誓。因为他们自己感觉到，所发誓的不过是虚无；更有甚者，他们还贿赂见证人，因为他竟敢公开藐视伪证的惩罚。关于这一点，你们中间谁没有伪证之罪？如今，指着神明发誓已无危险，因为指着\Person{凯撒}发誓会带来更重的顾忌，这本身便有损你们的神明；因为指着\Person{凯撒}起假誓的人，比违背对\Person{朱庇特}的誓言更容易受罚。\newline{}在轻视与嘲笑这两种相近的情感中，轻视较为体面，其傲慢带有某种荣耀；因为它有时源于自信、良心的安稳或天生的高尚心志。嘲笑则是一种更轻浮的情感，因此更直接指向苛责的傲慢。试想，你们将自己表现为多么伟大的嘲笑者！姑且不论你们在献祭时如何放纵这种情感——你们献上最贫瘠、最瘦弱的牺牲；或者从健康的动物中只取不能食用的部分，如头、蹄、拔下的羽毛和毛发，以及任何在家中被丢弃的东西。我也略过那些在粗俗和世俗之人看来构成你们祖先宗教的事物；然而，最有学问、最庄重的阶层（因为庄重与智慧在某种程度上自称源于学问），事实上总是最不敬畏你们的神明；即便他们的学识有所间断，也只是为了用更可耻的愚行和谎言弥补懈怠。我要从你们对那人的狂热喜爱说起——所有堕落的作家都从他那里获得梦境；你们对他赋予的尊荣，恰等于你们从神明那里夺去的，因你们抬高那嘲弄神明的人。我指的是\Person{荷马}。依我之见，他将神性的尊严贬低到人类状况的低级水平，用人的堕落和情欲浸染诸神；他使他们彼此争斗，胜负无常，如同成对的角斗士：他用凡人之手射出的箭伤及\Person{维纳斯}；他将\Person{玛尔斯}囚禁在锁链中十三个月，眼看就要死去；他让\Person{朱庇特}遭到一群天界叛徒的同样侮辱；或者为他因\Person{萨尔珀冬}而流泪；或者描绘他以最可耻的方式与\Person{朱诺}调情，通过描述和列举他的种种恋情来主张对她的乱伦欲望。自那以后，哪位诗人不曾以他们伟大君王为权威，诽谤诸神——要么泄露真相，要么捏造谎言？剧作家们，不论悲剧还是喜剧，岂不也曾让诸神成为他们戏剧中灾祸与报应的起因？姑且不提你们的哲学家，某种来自真理本身的灵感使他们奋起反抗诸神，并在他们骄傲的严厉与严格的戒律中免于恐惧。例如\Person{苏格拉底}。为轻视你们的神明，他指着橡树、狗和山羊发誓。虽然为此他被判处死刑，但\Person{雅典人}后来后悔了，甚至处死了他的控告者。通过这种行为，\Person{苏格拉底}的见证恢复了全部价值，而我得以用这句话回击你们：在他身上，你们认可了如今在我们身上被谴责的行为。除了这个例子，还有\Person{第欧根尼}——我不知道他嘲笑\Person{赫拉克勒斯}到何种程度；而罗马式的\Person{第欧根尼}——\Person{瓦罗}，向我们展示大约三百位约维斯，或应称作\Person{朱庇特}，全都没有头颅。你们其他轻浮的才智之士也用侮辱神明的方式为你们的享乐服务。仔细审视你们\Person{伦图路斯}和\Person{奥斯提乌斯}作品中渎神的美妙之处：他们在戏弄和玩笑中，到底是演员还是你们的神明成了笑柄？再者，你们多么愉快地接纳那些描述神明一切丑行的舞台文学！他们的威严在你们面前被某个不洁的身体玷污。某个神祇的面具，随你们的意愿，覆盖着某个可耻的卑微头颅。\Person{太阳}为他儿子的死亡而哀悼，伴随闪电，而你们却粗鲁地欢呼。\Person{库柏勒}为一位蔑视她的牧羊人叹息，而你们毫不脸红；你们平静地容忍那些歌颂\Person{朱庇特}风流韵事的歌曲。当然，当你们的神明以同样的热情在角斗士表演中舞蹈，为人血的流淌、为那些既是他们证明又是他们处决罪犯的污秽刑罚而欢欣，或者当你们的罪犯亲自扮演神明受罚时，你们拥有更虔诚的精神。我们常见在一个被肢解的罪犯身上，你们\Place{佩西努斯}的神\Person{阿提斯}；一个被活活烧死的可怜人扮演了\Person{赫拉克勒斯}。我们嘲笑你们正午神明游戏的乐趣，当\Person{冥父}、\Person{朱庇特}的亲兄弟，手持锤子拖走角斗士的残骸；当\Person{墨丘利}戴着有翼的帽子、拿着烧红的铁棒，用他的烙铁测试尸体是否真的无生命，还是假装死亡。谁能逐一追究这类细节？它们虽损毁神明的尊荣，令威严蒙羞，但都源于那些扮演者以及那些可以被如此描绘的人对神明的轻视。因此，我几乎不知道，你们的神明更有理由抱怨你们还是抱怨我们。一方面你们轻视他们，另一方面你们奉承他们；若对他们有任何失职，你们便用酬金平息；简言之，你们允许自己对神明随意而为。然而，我们却以一贯而完全的憎恶对待他们。\par

% structure: p0025 source=h2 target=subsection reason=relative-heading-rank; base=h2

\subsection{第十一章——关于驴头之荒谬妄言的驳斥}

在此事上，我们（据说）不仅犯有背弃公共宗教之罪，更引入了某种怪异的迷信；因为你们中有人臆想我们的天主是驴头——这是\Person{科尔内利乌斯·塔西陀}首先提出的一种谬论。他在其\WorkTitle{历史}第四卷中，论述犹太战争时，从该民族的起源开始描述，并就他们宗教的起源和名称发表了自己的见解。他叙述说，犹太人在旷野迁徙中，因缺水而困苦，借助一些野驴为向导才得以逃脱——他们以为这些野驴是在觅食牧草后前去寻水——因此犹太人对其中一头动物的形象加以崇拜。由此，我猜想，人们便推定我们，因与犹太宗教关系密切，也同样拜祭以同一象征形式所圣化的神物。然而，同一位\Person{科尔内利乌斯·塔西陀}——说实话，他在说谎方面最为饶舌——后来却忘了自己的说法，又叙述说\Person{大庞培}征服犹太人并攻占耶路撒冷后，进入圣殿，尽管仔细搜查了该处，却没有发现任何形象之物。那么，他们的天主应当在哪里呢？当然只能在这样一座令人难忘的圣殿里，它除了司铎外对所有的人严密关闭，陌生人不敢进入。但在此我应如何为自己即将说的话道歉呢？此刻我别无他图，只想泛泛地略作一二评论，这些评论同样适用于你们自己。那么，假设我们的天主是驴形人物，你们能否认在这件事上你们与我们有相同特征吗？（不只是他们的头，而是）整头驴，连同它们的守护神\OriginalTerm{厄波娜}{Epona}，确实都是你们崇拜的对象；而且你们将一切畜群、牛群和野兽，连同它们的厩棚都一并祝圣！也许这就是你们对我们的不满：当我们被各类牲畜崇拜者包围时，我们却独独崇拜驴子！\par

% structure: p0027 source=h2 target=subsection reason=relative-heading-rank; base=h2

\subsection{第十二章——崇拜十字架的指控。异教徒自身在神圣事物中大量使用十字架；不仅如此，他们的偶像本身也是以十字架框架形成的。}

至于那断言我们是\Quote{十字架的司祭职}的人，我们将视其为我们的同信仰者。十字架，就其材料而言，是木头的标记；而在你们中间，崇拜的对象也是木制偶像。只是，在你们那里偶像是人的形象，在我们这里木头本身就是其形象。目前暂且不论形状如何，只要材料相同：形式也无关紧要，如果它确实是神的真身。然而，如果在这方面出现差异的问题，那么请问，雅典的\Person{帕拉斯}或法罗斯的\Person{刻瑞斯}与制成十字架的木料之间有何区别，既然每一件都由粗糙的木料、无形之物、以及未成形木像的最初雏形来代表？每一根固定在地上的直立木料都是十字架的一部分，而且确实是其大部分体积。但一个完整的十字架归属于我们，当然有横梁和突出的座位。现在你们更无可推诿，因为你们只将一块残缺不全的木料奉献给宗教，而他人则将完整结构献为神圣用途。然而，实情终究是，你们的宗教\Emph{全是十字架}，正如我将要展示的。你们确实不知道，你们的神在起源上正是来自这可憎的十字架。现在，每一尊像，无论是木雕、石雕、金属铸造，还是用其他更贵重材料制成，都必然有塑造的手参与其形成。那么，这位塑造者，在做任何其他事之前，首先想到木十字架的形状，因为我们自己的身体也以自然姿势呈现出隐秘而隐藏的十字架轮廓。由于头向上抬起，背挺直，肩膀向两侧伸出，如果你只是将一个人双臂伸直平放，你就会做出十字架的大体轮廓。从这雏形和支柱出发，他涂上一层黏土，逐渐完成四肢，形成身体，并用他打算印在黏土上的形状覆盖内部的十字架；然后从这个设计，借助圆规和铅模，他准备好了一切，以便将其像制成大理石、黏土或他决定制作他的神的任何材料。（那么，过程就是：）在十字架形骨架之后是黏土；黏土之后是神。按照众所周知的程序，十字架通过黏土媒介转变为神。于是你们祝圣十字架，并且从它那里开始，祝圣的神灵获得其起源。举例来说，让我们以一棵树为例，它生长出枝干和树叶，并繁殖其同类，无论它是由橄榄仁、桃核、还是经过地下充分处理的胡椒粒所长出。现在，如果你移植它，或取其枝条用于另一植物，你会将繁殖的产物归于什么？难道不是归于谷粒、桃核或果仁吗？因为，正如第三阶段归于第二阶段，第二阶段同样归于第一阶段，所以第三阶段必须通过作为中间的第二阶段归于第一阶段。我们无需在此问题上继续讨论，因为根据自然法则，自然界每一种产物都将其生长追溯至其原始源头；正如产物包含在其最初原因中，该原因也与产物的特性相符。因此，既然在你们创造神的过程中，你们崇拜产生它们的十字架，那么这里就是原始的果仁和谷粒，你们的偶像崇拜的木料由此繁殖。例子不难找到。你们以宗教仪式将胜利奉为神明；胜利越带给你们欢乐，就越显尊贵。悬挂战利品的框架必须是十字架：这些可以说是你们游行庆典的核心。因此，在你们的胜利中，军营的宗教甚至使十字架成为崇拜对象；它崇拜你们的军旗，你们的军旗是其誓言的保证；它将你们的军旗置于\Person{朱庇特}本人之上。但所有那些像的炫耀和纯金的展示，不过是十字架的项链。同样，在你们的旗帜和军旗中——士兵们以同样神圣的谨慎守护着它们——你们有十字架的飘带和衣饰。我想，你们羞于崇拜朴素无华的十字架。\par

% structure: p0029 source=h2 target=subsection reason=relative-heading-rank; base=h2

\subsection{第十三章——崇拜太阳的指控被反驳。}

其他人，必须承认，更顾及良好的礼仪，认为太阳是基督徒的神，因为我们向东祈祷是众所周知的事实，或者因为我们使星期日成为庆祝日。那么怎样？你们做得少吗？你们中间很多人不是有时也假装崇拜天体，朝着日出的方向动嘴唇吗？总之，正是你们甚至把太阳纳入周历；你们选择它的日子，优先于前一日，视为一周中最适合完全禁浴、或推迟到晚上、或休息和宴饮的日子。采取这些习俗，你们故意偏离自己的宗教仪式，转向陌生人的仪式。因为犹太人在安息日和\Quote{洁净}的节期，以及灯烛的仪式、无酵饼的斋戒和\Quote{海滨祈祷}，所有这些制度和实践当然与你们的神无关。因此，为了让我从这个题外话回来，你们这些以太阳和星期日责备我们的人，应当考虑你们与我们的接近。我们离你们的\Person{萨图尔努斯}和你们的休息日并不远。\par

% structure: p0031 source=h2 target=subsection reason=relative-heading-rank; base=h2

\subsection{第十四章——关于奥诺科特斯的卑鄙诽谤：戴尔都良反问外邦人。}

传闻针对我们的天主引入了一项新的诽谤。不久之前，在你们那座城市中，一个最为堕落的恶徒，一个确实已抛弃自己宗教的人——实际上是一名犹太人，他只不过失去了自己的皮，当然是被野兽剥去的，他每日带着健全的身体为雇佣而进入斗兽场与之搏斗，从而处于即将失去皮的状态——在公众中携带了一幅我们的漫画，上面有这个标签：\OriginalTerm{奥诺科特斯}{Onocoetes}。这个形象有驴耳朵，身穿一件\OriginalTerm{托加袍}{toga}，手拿一本书，一只脚上有蹄子。民众竟相信了这个臭名昭著的犹太人。因为还有什么别的人群是一切反诽谤我们的种子床呢？因此，在整个城市中，奥诺科特斯成为人人谈论的话题。然而，既然它不过是\Quote{九日之奇}，因而不具有任何时间的权威，并且从其作者的性格来看也软弱无力，我将单纯以反驳的方式使用它来取悦自己。那么让我们看看，你们是否在这里也与我们为伍。现在，当我们的关注点在于畸形形象时，它们的形状如何并不重要。你们中间有狗头、狮头、牛角、公羊角和山羊角的神，羊形或蛇形，脚、头、背部长有翅膀。那么为什么如此醒目地标记我们独一的天主呢？你们自己中间就有许多奥诺科特斯。\par

% structure: p0033 source=h2 target=subsection reason=relative-heading-rank; base=h2

\subsection{第十五章——弑婴指控反诘异教徒}

既然我们在神祇方面处于同等地位，那么在祭献甚至崇拜方面，我们之间也没有区别——如果我可以从另一类证据来证实我们的比较。我们以杀死婴儿开始我们的宗教仪式或开启我们的奥迹。至于你们，既然你们自己关于人血和弑婴的行径已经从你们的记忆中淡去，我将在适当的地方适时提醒你们；我们现在推迟大部分例子，以免我们似乎处处在处理相同的话题。同时，正如我所说，我们之间的比较在另一观点下也并不失败。因为，如果我们在某种意义上是弑婴者，你们在任何其他意义上也很难被认为是如此；因为，虽然法律禁止你们杀死新生婴儿，但事情却是，没有任何法律被更逍遥地规避，而且是在公众的知情和整个时代的赞同下。然而，我们之间没有多大区别，只是你们杀死婴儿不是通过神圣的仪式，也不是（作为服务）献给天主。但是你们以更残忍的方式除掉他们，因为你们将他们遗弃在寒冷、饥饿和野兽之中，或者用更慢的淹死来摆脱他们。然而，如果在这方面我们之间确实存在任何不同，你们必须不忽视这样一个事实：你们扼杀的是你们自己的亲生孩子；这将在你们那方面补充，甚至大大加重，任何我们在其他方面所欠缺的。嗯，但我们据说享用了我们不敬的祭献！当我们推迟到更合适的地方讨论即使是这种做法在你们中间发现的任何相似之处时，我们在贪婪方面与你们没有太大区别。如果在一种情况中是不洁，而在我们的情况中是残忍，我们仍然处于同样的基础上（如果我可以承认我们的罪孽），在自然中，残忍总是与不洁相一致。但是，说到底，你们比我们做得少什么呢？或者，你们什么事不是做得比我们过分呢？我不知道，你们渴望人的内脏是否是一件小事，因为你们活生生地吞食成年人？真的，舔舐人血难道只是一件小事，当你们抽出本应活着的血？在你们看来，吃一个婴儿是一件轻松的事，当你们在它出生之前就整个消耗掉它？\par

% structure: p0035 source=h2 target=subsection reason=relative-heading-rank; base=h2

\subsection{第十六章——其他指控以同样方法驳斥。高贵罗马青年及其父母的故事}

我现在来到了熄灯、使用狗、并实践黑暗行为的时刻。在这一点上，我恐怕必须向你们屈服；因为我能对你们提出什么类似的指控呢？但是你们应该立即称赞我们的智慧，即我们以不伦之事显得端庄，因为我们设计了一个虚假的夜晚，以避免玷污真正的光明和黑暗，甚至认为摒弃地上的光、并戏弄我们自己的良心是合适的。因为凡是我们自己做的，当我们选择（怀疑）时，我们怀疑别人。至于你们的不伦之事，恰恰相反，人们在光天化日之下、或自然的黑夜中、或高天之下完全自由地享受它们；并且与它们的成功程度成正比，你们自己对结果的无知，因为你们在光天化日的完全认知下公开行淫。（然而，无知并不能将我们的行为隐藏于我们的眼睛，）因为在最黑暗之中，我们能够认出我们自己的恶行。你们非常清楚，根据\Person{克特西亚斯}的说法，波斯人生平与他们的母亲完全混杂，知道这个事实，毫无恐惧；而马其顿人众所周知，他们经常做同样的事，并且完全认可：因为有一次，当瞎眼的\Person{俄狄浦斯}出现在他们的舞台上时，他们以大笑和嘲笑的欢呼迎接他。演员惊恐地摘下他的面具，说：\Quote{先生们，我令你们不悦了吗？}\Quote{当然没有，}马其顿人回答，\Quote{你演得足够好；但是要么作者非常愚蠢，如果他编造了（这种自残作为对乱伦的赎罪），要么俄狄浦斯非常愚蠢，他那样惩罚自己；}然后他们彼此对喊道……\par

\SourceRecord{Translated by Peter Holmes.；Ante-Nicene Fathers；Vol. 3.；Edited by Alexander Roberts, James Donaldson, and A. Cleveland Coxe.；Buffalo, NY: Christian Literature Publishing Co.,；1885.；<http://www.newadvent.org/fathers/03061.htm>.}

% structure: p0001 source=h1 target=section reason=page-title

\section{外邦人（卷二）}

% structure: p0002 source=h2 target=subsection reason=relative-heading-rank; base=h2

\subsection{第一章：异教权威笔下的异教神。\Person{瓦罗}曾就此主题著书。他的三分法：本应固定不变者之善变性格。}

我们的辩护要求我们在此时与你们——可怜的外邦人——讨论你们众神的性格，甚至诉诸你们自己的良心，来判断它们究竟是如你们所愿的那样是真神，还是如你们不愿被证明的那样是假神。这是人类错误的实质部分，因着其作者的诡计，它从未摆脱对错误的无知，因此你们的罪责更大。你们的眼睛睁开，却看不见；耳朵畅通，却听不见；虽然心在跳动，却仍然迟钝，理智也未能理解它所认知之事。如果（你们崇拜的）巨大乖谬能因一个单一的反对而被打破，我们便已准备好一个异议，即我们已知你们所有的神都是人设立的，这一事实本身便摧毁了对真天主的一切信仰；因为，当然，任何曾经有开端的事物都不能正确地被视为神圣。但事实是，有许多事情使良心的柔弱在故意错误的麻木中变得刚硬。真理被（敌人的）庞大力量所围困，然而在其内在的力量中她是多么安全！很自然，即使从她的仇敌中，她也能赢得她所愿意的人作为朋友和保护者，并击倒所有攻击者的大军。因此，我们的斗争正是针对这些事物——反对我们祖先的制度、反对传统的权威、统治者的法律、智者的推理；反对古旧、习俗、服从；反对先例、奇事、奇迹——所有这一切都在巩固你们众神的虚假体系中。因此，我们愿意逐步跟随你们自己的注疏——这些注疏是你们从各种神学中提取出来的（因为在这类事情上，学者的权威对你们来说比事实的见证更有分量）——我选取并缩写了的\Person{瓦罗}作品；因为他在他的著作《论神圣之事》中，从古代的摘要里搜集材料，为我们展示了一个有用的指南。现在，如果我向他询问谁是众神的精巧发明者，他指向哲学家、各民族或诗人。因为他在分类众神时做了三重区分：一类是\Emph{物理的}，由哲学家处理；另一类是\Emph{神话的}，由诗人不断重复；第三类是\Emph{民族的}，各民族各自采用。因此，当哲学家凭自己的猜想精巧地构成他们的物理（神学），当诗人从寓言中得出他们的神话（神学），而（各个）民族凭自己的意愿锻造他们的民族（多神教）时，真理究竟在哪里？在猜想中？好吧，但这只是一种不确定的概念。在寓言中？但它们充其量是荒谬的故事。在流行说法中？然而，这种看法不过是混杂和地域性的。哲学家那里的一切都是不确定的，因为他们的变化多端；诗人那里的一切都是无价值的，因为不道德；各民族那里的一切都是不规则和混乱的，因为取决于他们的纯粹选择。然而，天主的本性——如果你们所关注的是真实的——具有如此确定的特征，以至于不能从不确定的思辨中得出，不能受无价值的寓言玷污，也不能由混杂的臆想决定。它确实应当被视为——正如它实际的那样——确定的、完整的、普遍的，因为它确实是属于万有的。那么，我该相信哪一位神呢？一个被模糊怀疑所衡量过的？一个被历史所揭露的？一个被社群所发明的？如果我根本不信任何神，也比信一个可疑的、或充满羞耻的、或任意选择的神要好得多。\par

% structure: p0004 source=h2 target=subsection reason=relative-heading-rank; base=h2

\subsection{第二章：哲学家未能发现天主。其思辨之不确定与混乱。}

但是，自然哲学家的权威被\Emph{在你们中间}当作智慧的专有财产来维护。\Emph{你们意指}的当然是哲学家们那种纯粹而简单的智慧，它主要通过那种源于对真理无知的观点多样性来证明自身的软弱。现在，有哪个聪明人会如此缺乏真理，以至于不知道天主是智慧本身和真理的圣父和主？此外，还有撒罗满所传的那道神谕：\Quote{敬畏主，}他说，\Quote{是智慧的开端。}但畏惧源于知识；因为一个人如何能畏惧他一无所知的东西呢？因此，凡拥有敬畏天主的人，即使对其他一切无知，若达到了对天主的认识与真理，也将拥有圆满而完美的智慧。然而，这正是哲学所未能清楚领悟的。因为，尽管哲学家们以其好探究的性格去钻研各种学问，似乎出于圣经的古远而研究了圣经本身，并从中得出了一些观点；但因他们窜改了\Emph{这些推论}，便证明他们要么完全轻视了圣经，要么并未全心相信它们，因为在其他情况下，真理的纯朴也被不纯正信仰的过度拘谨所动摇，因此他们随着荣誉欲望的增长，把它们改造成了自身心智的产物。其结果是，即便是他们所发现的东西也退化成了不确定，从一两滴真理中涌出了滔滔不绝的争论。因为他们在简单发现天主之后，并未按其所发现的那样阐述祂，反而争论祂的属性、祂的本性，甚至祂的居所。柏拉图学派确实认为祂关心世间事物，既是其安排者又是其审判者。伊壁鸠鲁学派把祂\Emph{视为}冷漠、惰怠，可以说是一个虚无。斯多亚学派相信祂在世界之外；柏拉图学派则认为祂在世界之内。他们如此不完美地承认的那位天主，他们既不能认识也不能畏惧；因此他们不能成为智者，因为他们确实偏离了智慧的开端，即\Quote{敬畏天主}。不乏证据表明，在哲学家当中，不仅存在无知，而且对神性实际抱有怀疑。第欧根尼被问及天上发生何事时，回答说：\Quote{我从未上去过。}又问是否有任何神祇，他答道：\Quote{我不知道；只应该存在神祇。}当克罗伊斯询问米利都的泰勒斯对诸神的看法时，后者思考了片刻，回答说：\Quote{没有。}甚至苏格拉底也以一种确定的口吻否认了你们那些神祇。然而他却以同样的确定要求向埃斯库拉皮乌斯献上一只公鸡。因此，当哲学在为天主下定义的实践中被发现如此不确定且前后矛盾时，它对于自己无法清楚确定的那一位能有怎样的\Quote{畏惧}呢？我们被教导要相信世界是神。因为自然神学家这一类人得出结论说它是这样，他们传下了关于诸神的观点，以至于斯多亚派的狄奥尼修斯将它们分为三类。他认为，第一类包括那些最明显的诸神，如太阳、月亮\Emph{和}星辰；第二类是那些不明显的，如尼普顿；最后一类是从人类状态过渡到神性的那些，如赫拉克勒斯\Emph{和}安菲阿拉俄斯。同样，阿塞西劳斯将神性分为三种形式——奥林匹亚、星辰、泰坦——出自克卢斯与特拉；由此经过萨图尔努斯与俄普斯产生了尼普顿、朱庇特与奥尔库斯，以及他们的全部后代。学园派的色诺克拉底做了二分——奥林匹亚与泰坦，二者出自克卢斯与特拉。大多数埃及人相信有四位神祇——太阳、月亮、天和地。德谟克利特推测诸神来自天上的火。芝诺也认为其本性类似于此。因此瓦罗也说火是世界的灵魂，在世界中火统治万物，正如灵魂在我们体内所做的那样。但这一切极其荒谬。因为他说，灵魂在我们体内时，我们存在；而一旦它离开我们，我们就死亡。因此，当火以闪电离开世界时，世界就终结了。\par

% structure: p0006 source=h2 target=subsection reason=relative-heading-rank; base=h2

\subsection{第三章 自然哲学家主张元素的神性；该学说的荒谬性被揭露。}

从这些意见的演变中，我们看到你们的自然派\Emph{哲学家}被迫主张元素是神，因为他们声称其他神是从元素中产生的；因为只有从神才能生出神。现在，虽然我们应在适当的地方，即在诗人的神话部分更充分地考察这些其他神，但既然我们必须同时就他们与当前类别的联系加以讨论，我们或许甚至从他们的当前类别，一旦我们转向诸神本身，就能成功地证明：那些据说是从元素中产生的神，绝不可能显得是神；由此我们立即有一个推定：元素不是神，因为从元素所生的不是神。同样，当我们表明元素不是神时，我们将按照自然关系的法则得到一个推定性论证：那些其父母（在此即元素）不是神的东西，不能被正确地主张为神。这是一个确定之点：神由神而生，缺乏神性者由非神圣者所生。现在，就你们的哲学家所论述的世界而言（因为我将\Emph{宇宙}这个术语用于最广泛意义上的世界），它包含元素，并将元素作为其组成部分来服务（因为无论其自身状态如何，其元素和各组成部分的状况当然也是如此），它必然要么由某个存在者所形成，如\Person{柏拉图}的卓见，要么不由任何存在者所形成，如\Person{伊壁鸠鲁}的严酷意见；而既然它是被形成的，因而有开始，它也必然有终结。因此，那在开始之前不存在、在终结之后也将不再存在的东西，当然绝不可能显得是神，因为它缺少神性的本质特征——永恒，而永恒被认为既无始也无终。然而，如果它根本不是被形成的，因此应被视为神圣——因为作为神圣者，它本身既无开始也无终结——那么，何以有些人将生成归于元素，认为它们是神，而斯多葛派却否认任何事物可以从神而生？同样，何以他们希望那些他们设想为由元素所生的东西被视为神，而他们却否认神可以被生？现在，凡适用于宇宙的，也必须适用于元素，我是指天、地、星辰和火，\Person{瓦罗}徒然地提议你们相信这些是神，而且是诸神的父母，这与他所宣称的神不可能有生成和出生相矛盾。这位\Person{瓦罗}曾证明大地和星辰是有生命的。但如果是这样，按照有生命之物的本性，它们也必然是可朽的；因为虽然灵魂显然是不朽的，但这一属性仅限于灵魂本身，并不延伸至与灵魂结合的东西，即身体。然而，没有人会否认元素有身体，因为我们既能触摸它们，也被它们触摸，并且我们看到某些物体从它们那里落下。因此，如果它们是有生命的，撇开灵魂的原则（就它们作为身体的条件而言），它们就是可朽的——当然不是不朽的。然而，何以元素在\Person{瓦罗}看来是有生命的？因为，确实，元素有运动。然后，为了预先反驳可能来自对方的主张——即许多其他东西也有运动，如车轮、马车、其他若干机器——他主动声明，他只相信那些自身运动而无任何外在明显推动者或驱动者的东西是有生命的，如同车轮的明显推动者、马车的推进器或机器的操纵者那样。那么，如果它们没有生命，它们就没有自身的运动。现在，当他如此断言一种不明显的动力时，他指出的是他应当寻求的东西，即运动的创造者和控制者；因为我们并不能因为看不见某物就相信它不存在。相反，我们必须更深入地探究看不见的东西，以便更好地理解显见之物的性质。此外，如果（你们承认）只有那些显现并且仅仅因为显现而被认为存在的东西才存在，那么你们又何以承认那些不显现的东西是神？再者，如果那些没有存在的东西似乎存在，那么那些似乎没有存在的东西为什么不能存在？例如，天上的众星体的推动者。那么，就算事物有生命是因为它们自身运动，并且它们自身运动时并非由他物推动：这仍然不意味着它们必须立刻就是神，因为它们是活物，甚至也不因为它们自身运动；否则，有什么能阻止所有动物——它们都自身运动——被算作神呢？这当然被埃及人容许，但他们的\Emph{迷信}虚荣另有根据。\par

% structure: p0008 source=h2 target=subsection reason=relative-heading-rank; base=h2

\subsection{第四章 \OriginalTerm{神}{\GreekText{Θεός}} 一词的错误衍义。此名指示真天主。\OriginalTerm{天主}{God} 无形体且非物质。\Person{泰勒斯}轶事。}

有些人断言，诸神（即\OriginalTerm{诸神}{\GreekText{θεοί}}）之所以如此称呼，是因为动词\OriginalTerm{奔跑}{\GreekText{θέειν}}和\OriginalTerm{被移动}{\GreekText{σείσθαι}}意指\Emph{奔跑}和\Emph{被移动}。因此，这个术语并不表示任何威严，因为它源自奔跑和运动，而非源于任何神性的统治权。然而，我们所敬拜的至高天主也被称为\OriginalTerm{天主}{\GreekText{Θεός}}，但在他内并无任何\Emph{行进}或\Emph{运动}的迹象，因为他是任何人所不能见的，显然，这个词必然有某种其他的来源，并且他内在本有的神性特质必定已被发现。因此，撇开那种巧妙的解释，更可能的是，诸神并非因\Emph{奔跑}和\Emph{运动}而被称作\OriginalTerm{诸神}{\GreekText{θεοί}}，而是这个术语借用了真天主的称号；因此，你们将\OriginalTerm{诸神}{\GreekText{θεοί}}这个名号给予了那些你们同样为自己捏造的诸神。如今，此事有一个明显的证明，即你们实际上将共同的称谓\OriginalTerm{诸神}{\GreekText{θεοί}}给予了你们所有的那些神，在他们身上并没有\Emph{行进}或\Emph{运动}的属性被指示出来。因此，当你们同样轻易地将他们称为\OriginalTerm{诸神}{\GreekText{θεοί}}和\Emph{不可移动}时，这既偏离了该词的含义，也偏离了神性的观念，若以\Emph{行进}和\Emph{运动}的概念来衡量，神性便被搁置了。但如果那\Emph{神圣}的名号特别指代神性，且对于\Emph{真}天主而言是单纯真实的、并非牵强附会的解释，却被以借用的意义转移到了你们选择称之为神的其他对象上，那么你们应当向我们证明，在他们之间也存在本质上的共同性，以便他们共同的称号可以正确地取决于他们的本质合一。但是，真天主仅仅因为不是感官对象，就无法与那些可被视觉和感觉（感觉就足够了）认知的虚假神祇相比较；因为这等于清楚地说明了隐晦证明与显明证明之间的区别。既然元素对所有人都是显而易见的，而相反，天主却是无一人可见的，你们怎能从你们未曾看见的那一部分，转而对你们所见之物做出判断呢？因此，既然你们没有必要在感知或理性中将它们结合，为什么你们又在名称上结合它们，目的是也在能力上结合它们呢？请看，就连\Person{芝诺}也将世界的物质与天主分开：他说后者渗透前者，如同蜂蜜穿过蜂巢。因此，天主与物质是两个词，也是两样东西。事物之间的差异与词之间的差异成正比；物质的状态也遵循其称谓。既然物质不是天主，因为其称谓本身就教导我们如此，那么那些内在于物质中的事物——即元素——又怎能被视为诸神呢？因为组成成分不可能与整体不同质。但我与那些物理学的奇思异想有何干系？人的心智最好上升到世界状态之上，而不是屈尊于不确定的思辨。\Person{柏拉图}认为世界的形状是圆的。我想，他用圆规将别人所设想的方形、有角形状弄圆了，因为他苦心要让人们相信世界根本没有开端。然而，\Person{伊壁鸠鲁}曾说过：\Quote{在我们之上的东西与我们无关，}却仍想窥视天空，并发现太阳的直径是一英尺。至此你必须承认，人们在涉及天体对象时也是吝啬的。随着时间的推移，他们雄心勃勃的概念进步了，太阳也随之增大了它的圆面。于是，逍遥学派将其标示为一个更大的世界。现在，请告诉我，这种对臆测性思辨的渴望有什么智慧？尽管其断言充满自信，但那种无用的、刻意的、吹毛求疵的好奇心，以巧妙的语言装饰，又能向我们提供什么证据呢？因此，\Person{米利都的泰勒斯}完全活该，当他走路时用尽所有眼睛观星，却不幸掉进井里，并被一个埃及人无情地嘲笑，那人对他说：\Quote{难道是因为在地上找不到可看的东西，你就认为应该把目光局限在天上吗？}因此，他的跌落是哲学家的一个象征性图景；我指的是那些坚持将研究用于虚荣目的的人，因为他们沉醉于对自然对象的愚蠢好奇，而他们本应明智地将目光转向他们的创造者和治理者。\par

% structure: p0010 source=h2 target=subsection reason=relative-heading-rank; base=h2

\subsection{第五章 物理理论的延续。进一步反对元素神性的理由。}

那么，我们为何不采纳那远为合理的见解呢？它显然源自人的常识和纯朴的推理。连\Person{Varro}也考虑到这一点，他说诸元素被认为是神圣的，因为离开它们的协同，任何事物都无法产生、滋养或用于维持人类生命与大地，因为就连我们的身体和灵魂本身，若无元素的调节，也无法自足。正是藉着元素，世界才普遍适于居住——这一结果是通过按地带分布的和谐协调而实现的，除非人类居住因严寒或酷热而无法实现。为此，人们将以下事物视为神：太阳，因为它自本身发出白昼之光，以其温暖使果实成熟，并以固定的周期度量岁月；月亮，它是黑夜的慰藉，以其治理掌管月份；还有星辰，它们是农事耕作中须观察的某些季节的明确标志；最后，还有天空（在其下）、大地（在其上）以及中间的空间（在其中），万物都为人的益处协力合作。不仅是由于它们的恩惠影响，人们认为神圣性可与元素并存，而且由于它们的相反特质，例如那些通常可称之为\Quote{它们的愤怒与恼怒}所产生的事物——如雷、冰雹、干旱、致病的风、洪水、地裂和地震：这些都相当合理地被视为神，无论其性质是因有益而受敬仰，还是因可怕而受畏惧——它们实为助益与伤害的至高施与者。但在社会生活的实际实践中，人的行为与感受是这样的：他们并不对产生帮助或伤害的事物本身表达感谢或责备，而是对藉其力量和权能实现这些事物运作的那一位。因为即使在你们的娱乐中，你们也不会将花冠奖给长笛或竖琴，而是奖给以其令人愉悦的技巧的力量驾驭那长笛或竖琴的乐师。同样，当一个人身体不适时，你们不会感谢毛织物、药物或膏药，而是感谢那些以其照料和智慧使治疗生效的医生。同样，在不幸事件中，被剑刺伤的人不会将伤害归咎于剑或矛，而是归咎于敌人或强盗；被倒塌房屋掩埋的人不会责备瓦片或石头，而是责备建筑的陈旧；同样，遇难的水手不会将灾祸归咎于礁石和波浪，而是归咎于风暴。这也是合理的；因为可以肯定，所发生的每件事都必须归因于其发生的工具，而是归因于藉其而发生的那一位；因为他是事件的首要原因，他指定事件本身以及藉以发生之工具（因为在所有事物中都有这三个特定要素——事件本身、其工具和其原因），因为他自己，那愿意某事发生的人，在他所愿意的事物或事件藉以发生的工具之前就已显现。因此，在其他一切场合，你们的行为是正确的，因为你们考虑作者；但在自然现象中，你们的规则反对那种自然原则——该原则促使你们在所有其他情况下作出明智判断——你们将作者的最高地位隐去，反而考虑所发生的事物，而非其发生者。于是，你们便以为权力与统治属于元素，而它们不过是奴仆和职能者。现在，我们在内中追溯一位工匠和主人时，岂不正是根据你们赋予权能的那些元素的既定职能，揭露了它们奴役状态的艺术结构吗？但神不是奴仆；因此，凡具奴仆性质者都不是神。否则，他们应向我们证明，按事物的通常进程，自由由无纪律的放纵促进，专制由自由促进，而神的权能意味着专制。因为如果天体在特定轨道、固定季节、适当距离\Emph{并且}相等间隔内完成其运行——这是为时间周期及其引导而按法律设立的——那么，从它们条件的遵守和运作的忠实，岂不能得出这样的结果：由于你们注意到其轨道运行的重复和变化的精确性，并记住其重复是何等持续不断，你们就会相信有一位主宰的权能统御它们，整个世界的管理都服从于它，甚至达于人类的使用与损害？因为你们不能假装这些现象只为自己行动和关怀，而不为人类利益贡献什么，因为你们主张元素是神圣的，唯一的理由是你们从它们那里体验到对自己的益处或损害。因为如果它们只造福自己，你们就无须感谢它们。\par

% structure: p0012 source=h2 target=subsection reason=relative-heading-rank; base=h2

\subsection{第六章 天体的变化：证明它们不是神圣的。从自然神向神话神的过渡。}

来吧，你们是否承认，天主在其运行中不仅毫无奴役性，而且以无缺的完整性存在，不应被减弱、中止或毁灭？那么，如果祂曾受变化，祂的一切福乐就会消失。然而看看那些星体；它们都经历变化，并清楚地证明了这一事实。月亮告诉我们它的亏损有多大，当它恢复全形时；你们已经习惯于在水镜中测量它更大的亏损；因此我不再需要相信占星家所断言的一切。太阳也时常经受日蚀的考验。尽管你们可以尽力解释这些天体偶然事件的方式，但天主既不能变小，也不能消失。因此，那些人类学问的支撑是徒劳的，它们通过巧妙地编织猜想，歪曲了智慧和真理。此外，按照你们自然的思维方式，确实发生这样的情况：说得最好的人被认为是最真实的，而说真话的人却被认为说得最好。现在，凡是仔细考察事物的人，必定会承认，我们一直在讨论的这些元素更有可能是受某种规则和指导的，而不是它们有自身的运动，并且由于受到管辖，它们不能是神。然而，如果有人在这一点上犯了错误，那么像你们自然哲学家那样简单地犯错比思辨地犯错更好。但与此同时，如果你们考虑\Emph{神话}学派的特点（并将它与\Emph{自然}学派比较），我们已经看到脆弱的人在后者中所犯的错误实际上是更可敬的，因为它将神性归于那些它认为在感觉上属于\Emph{超人}的东西，无论是就其位置、能力、大小还是神性而言。因为你们认为高于人的东西，你们相信非常接近天主。\par

% structure: p0014 source=h2 target=subsection reason=relative-heading-rank; base=h2

\subsection{第七章 神话一类的诸神。诗人们在这样的事上是非常不可靠的权威。\Person{Homer}与神话诗人。为何不虔诚。}

然而，若要论及归诸诗人的\Emph{神话的}神祇类别，我几乎不知是该仅试图将它们与我们自己的\Emph{凡俗}平庸等量齐观，还是必须用神性的证据宣称它们为神，如同非洲的\Person{莫普索斯}和玻俄提亚的\Person{安菲阿劳斯}那样。\newline{}现在我确实必须只略微触及此类，其更全面的考察将在适当的地方进行。\newline{}同时，这些人只是凡人，从你们并不一贯称他们为神而称他们为英雄的事实来看，是清楚的。\newline{}那么为何要讨论这一点呢？尽管必须将神性的尊荣归于死者，但当然不是归于他们本身如此。\newline{}看看你们自己的做法：当你们以同样的狂妄过度用你们的君王的坟墓玷污天堂时，你们所尊崇并赋予神化的福乐的不正是那些因正义、德行、虔诚和所有此类卓越而显赫的人吗？你们甚至甘愿为了这些人物而发假誓招致轻蔑也心满意足。\newline{}另一方面，你们岂不剥夺不敬和可耻之人甚至旧有的人类荣耀的奖赏，撕毁他们的法令和头衔，推倒他们的雕像，并污损他们铸币上的肖像吗？\newline{}然而，那洞察万物、赞许并奖励善者的祂，岂会将自己取之不尽的恩宠和慈悲的德能公然玷污呢？\newline{}难道人就被允许拥有一座特别的谨慎和公义之山，以便他们能明智地选择和倍增自己的神祇吗？\newline{}服侍君王和王侯的侍从岂能比服侍至上天主的人更纯洁吗？\newline{}你们确实恐怖地背弃流放者和被放逐者、贫穷和软弱者、出身微贱和卑劣生活者；然而你们竟以法律制裁来尊崇不贞洁者、通奸者、强盗和弑亲者。\newline{}那些不配为人的角色却被相信为神，我们应视此为笑料还是义愤之事呢？\newline{}再者，在你们这诗人们所颂扬的神话类别中，你们在良心的纯洁及其持守上的行为是何等摇摆不定！\newline{}因为每当我们指责你们诸神的悲惨、可耻和凶恶的实例时，你们便以假借诗意的\OriginalTerm{许可}{licence}为由，辩称它们只是寓言；\newline{}每当我们主动对所说的这种诗意的\Emph{许可}表示沉默的轻蔑时，你们不仅毫无恐惧，甚至进而表示尊敬，并视之为不可或缺的（高雅的）艺术之一；不仅如此，你们还借此进行高等阶层的研习，作为你们文学的基石。\newline{}\Person{柏拉图}认为诗人应当被放逐，因为他们是诸神的诽谤者；他甚至要把\Person{荷马}本人从共和国中驱逐出去，尽管如你们所知，他是所有诗人中加冕的魁首。\newline{}但既然你们这样接纳并保留他们，当他们对你们的诸神透露这些事情时，你们为何不相信他们呢？\newline{}如果你们相信你们的诗人，你们又怎么会敬拜（如他们所描述的）这样的神呢？如果你们敬拜他们仅仅是因为你们不相信诗人，你们又为何赞扬这些撒谎的作者，毫无畏惧得罪那些你们尊崇其诽谤者的人物呢？\newline{}当然，诗人身上不能期望对真理的尊重。\newline{}但当你们说他们只是把人死后变成神时，你们岂不承认死前这些所谓的神祇不过是人吗？\newline{}那么，那些曾是人者易遭受人类灾祸、罪行或寓言的不名誉，这有什么奇怪的呢？\newline{}事实上，当你们依据他们的狂诗在某些情况下安排你们的礼仪时，你们岂不正是相信你们的诗人吗？\newline{}为什么\Person{刻瑞斯}的女司祭被强暴，如果不是因为\Person{刻瑞斯}遭受了类似的暴行呢？\newline{}为什么别人的孩子被献祭给\Person{萨图尔努斯}，如果不是因为他没有饶恕自己的孩子呢？\newline{}为什么一个男性为了纪念伊达山的女神\Emph{库柏勒}而受阉割，如果不是因为那个（不幸的）青年对她过于轻蔑，因她恼怒他竟敢违抗她的爱而被阉割呢？\newline{}为什么\Person{赫拉克勒斯}对\Place{拉努维乌姆}的贵妇们不是\Quote{美味佳肴}，如果不是因为妇女们最初对他犯下的冒犯呢？\newline{}诗人无疑是撒谎者。然而，这既不是因为\Emph{他们告诉我们}你们的诸神作为人时做了这些事，也不是因为他们将神性的丑闻归之于神性状态，因为在你们看来，有神存在（即使不是这种品性）似乎比有这种品性（即使不是神）更可信。\par

% structure: p0016 source=h2 target=subsection reason=relative-heading-rank; base=h2

\subsection{第八章 不同民族的诸神。\Person{瓦罗}的\OriginalTerm{外邦人的类别}{Gentile Class}。它们的低劣。这种\OriginalTerm{歪曲的神学}{Perverse Theology}大部分取自圣经。\Person{塞拉比斯}是对\Person{若瑟}的歪曲。}

剩下的是分散在各民族中的\Emph{外邦}神祇一类；这些神祇纯属任意杜撰，并非出于对真理的认识；我们对它们的了解来自\Emph{不同种族}的私人观念。我想，天主是无所不在、处处临在、全能的神——是所有人都应当敬拜、所有人都应当事奉的对象。既然连全世界共同敬拜的那些神祇，尚且缺乏其真神性的证据，那么那些连其敬拜者都未能真正发现的神祇，岂不更应如此？因为一种如此残缺不全、完全不为世人所知的神学，又能有什么可信的权威呢？有多少人曾见过或听说过叙利亚的\Person{阿塔尔加提斯}、非洲的\Person{科莱斯蒂斯}、摩尔人的\Person{瓦尔苏蒂娜}、阿拉伯的\Person{奥博达斯}和\Person{杜萨里斯}，或诺里克的\Person{贝莱努斯}，以及\Person{瓦罗}所提及的——卡西努姆的\Person{德卢恩蒂努斯}、纳尔尼亚的\Person{维西迪亚努斯}、阿蒂纳的\Person{努米特尔努斯}，\Emph{或}阿斯库卢姆的\Person{安卡里亚}？又有谁清楚知道武尔西尼的\Person{诺尔提亚}？即使她们的名字，除去加于其上的凡人姓氏之外，其本身的价值也并无区别。我常常嘲笑各个市镇中那些小圈子的神祇，它们的尊荣仅限于本城的城墙之内。这种任意杜撰神祇的放荡行为已经发展到何种程度，埃及人的迷信做法向我们表明；因为他们甚至敬拜自己本地的动物，\Emph{例如}猫、鳄鱼和蛇。因此，他们还将一个人神化，这算不得什么大事——我指的是那个人，不仅埃及或希腊，而是整个世界都在敬拜，非洲人也凭着他起誓；关于他的境况，那些有助于我们推测、并赋予我们的知识以真理外观的一切，都记载在我们的（圣经）文献中。因为你们的那个\Person{塞拉皮斯}原本是我们的一位圣人，名叫\Person{若瑟}。他是兄弟中最年幼的，但智力超群，因被嫉妒而被卖到埃及，成了该国法郎王宫廷中的奴隶。遭到不贞洁的王后纠缠，当他拒绝顺从她的欲望时，她便反过来诬告他，向法郎禀报，结果他被投入监牢。在那里，他通过正确解释一些（同囚者的）梦境，展示了其神圣灵感的威力。与此同时，法郎也做了些可怕的梦。若瑟应召被带到法郎面前，得以解释这些梦。他讲述了在监牢中证明他正确解梦的证据，然后向法郎揭示了他的梦：那七头肥美健壮的牛象征七个丰年；同样，七头瘦弱的牛预示着随后七年的饥荒。因此，他建议利用前七年的丰盈为未来的饥荒做好防备。法郎信了他。后来发生的一切表明他是何等智慧、何等始终圣洁，以及如今何等必要。于是法郎派他治理全埃及，以便为埃及储备粮食，并从此治理其国政。他们称他为\Person{塞拉皮斯}，源于装饰他头部的头巾。这头巾的如配克状形状标志着他储备粮食的记忆；而头巾边缘装饰的麦穗则证明粮食供应的管理全在他头上。出于同样的原因，他们还造了一尊神圣的狗像，他们认为这狗是在冥界（充当哨兵），并将其放在他的右手下，因为埃及人的关切都集中在他的手下。他们在他身边放置了\Person{法里亚}，她的名字表明她是法郎的女儿。因为在他所有其他恩赐和赏报之外，法郎还将自己的女儿嫁给了他。然而，既然他们已经开始敬拜野兽和人类，他们便将两者结合为一种形象，即\Person{阿努比斯}，其中恐怕更可以清楚地看出一个与自身交战、抗拒其君王、被外邦人鄙视、甚至拥有奴隶的食欲和狗之污秽本性的民族所尊奉的、关于其自身品性和状况的证据。\par

% structure: p0018 source=h2 target=subsection reason=relative-heading-rank; base=h2

\subsection{第9章 罗马的权力。所有异教神话的罗马化面貌。\Person{瓦罗}的三重划分受到批评。罗马英雄（包括\Person{埃涅阿斯}）受到不利的审视。}

以上便是我们不得不提及的、与\OriginalTerm{瓦罗}{Varro}的三分诸神分类相关的一些较明显或较引人注目之处，以便就自然类、神话类和民族类诸神给出充分的回答。然而，既然如今取代（异教崇拜为真宗教）的责任已不再属于哲学家、诗人或民族——尽管他们传递了迷信——而是属于统治的罗马人，他们接受了传统并赋予了它广泛的权威，那么我们现在必须面对人类广泛错误的另一面向；不，必须用\Emph{我们的斧头}砍倒另一片森林，这片森林用从四面八方收集的迷信种子遮蔽了堕落崇拜的初始阶段。好吧，就连罗马人的诸神也从（同一位）\Person{瓦罗}那里得到了三分分类：\Emph{确定的}、\Emph{不确定的}和\Emph{精选的}。何等荒谬！当他们已经拥有确定的神时，何需不确定的神？除非他们确实想效仿雅典人的愚蠢；因为在雅典有一座祭坛上刻有铭文：\Quote{\OriginalTerm{致未知之神}{To the unknown gods}}。那么，人会崇拜他毫不知晓的东西吗？再者，既然他们有确定的神，他们本应满足于此，无需精选的神。这种需求甚至表明他们是渎神的！因为如果神明像洋葱一样被挑选，那么未被选中的就被宣布为毫无价值。现在我们这边承认罗马人有两类神：\Emph{共同}的和\Emph{特有的}；换言之，即他们与其他民族共有的神，以及他们自己创造的神。这些不是被称为\Emph{公共}神和\Emph{外来}神吗？他们的祭坛告诉我们：在卡尔娜（Carna）的神庙里有（一个）外来神的例子，在帕拉蒂尼（Palatium）则有公共神的例子。既然他们的共同神已经被包含在自然类和神话类中，关于它们我们已经说得够多了。我想谈谈他们特有的神祇种类。那么，我们应该钦佩罗马人那第三类神——\Emph{他们敌人的神祇}，因为没有其他民族为自己发现如此大量的迷信。他们的其他神祇我们分为两类：那些从人类变成神的，以及那些以其他方式起源的。既然对于死者神化的相同貌似合理的借口是他们的生活有功绩，我们不得不提出相同的反驳：他们中没有一个人值得如此费力。他们相信的慈父\OriginalTerm{埃涅阿斯}{Aeneas}从未荣耀过，被一块石头——一种用来掷狗的粗俗武器——击倒，所受的伤同样不光彩！但这个\Person{埃涅阿斯}结果成了国家的叛徒；是的，几乎与\OriginalTerm{安忒诺尔}{Antenor}一样。如果他们不相信这对他属实，他至少在国家燃烧时抛弃了同伴，必然不如那位迦太基妇女，她的丈夫\OriginalTerm{哈斯德鲁巴}{Hasdrubal}像我们\Person{埃涅阿斯}一样以温和的怯懦向敌人恳求时，她拒绝跟随他，而是带着孩子冲入燃烧的迦太基的火焰，仿佛投入她（亲爱但毁灭的）国家的怀抱。他是\Quote{虔诚的埃涅阿斯}吗？因（救出）他年幼的独子和年迈衰弱的父亲，却抛弃了\OriginalTerm{普里阿摩斯}{Priam}和\OriginalTerm{阿斯提阿那克斯}{Astyanax}？罗马人更应憎恨他；因为他们为了保卫自己的君主和王室，甚至交出孩子、妻子和最珍贵的抵押品。他们将\OriginalTerm{维纳斯}{Venus}之子神化，而且这是在丈夫\OriginalTerm{武尔坎}{Vulcan}完全知晓并同意的情况下，甚至没有\OriginalTerm{朱诺}{Juno}的反对。如果儿子因对父母的孝心而拥有天上的席位，那么为什么\OriginalTerm{阿尔戈斯}{Argos}那些高贵的青年不被视为神？因为他们为了在举行某些神圣仪式时拯救母亲脱离罪责，以超人的虔诚，将自己套上她的车，将她拖到庙里。为什么那位从自己乳房哺育了在狱中挨饿的父亲、以非凡孝心的女儿不被封为女神？关于\Person{埃涅阿斯}还能讲述什么其他光辉成就，除了在\OriginalTerm{劳伦图姆}{Laurentum}战场上他无处可见？也许他依其本性，第二次从战斗中逃走。类似地，\OriginalTerm{罗慕路斯}{Romulus}死后成了神。是因为他建立了城市？那么其他人呢？那些建造城市的人，甚至包括妇女？当然，\Person{罗慕路斯}还杀死了他的兄弟，并狡猾地掳掠了一些外邦少女。因此他当然成了神，也因此成了\OriginalTerm{奎里努斯}{Quirinus}（\Quote{矛之神}），因为那时他们的父亲不得不为他使用长矛。\OriginalTerm{斯特库鲁斯}{Sterculus}做了什么事值得神化？如果他辛勤地用\OriginalTerm{粪肥}{stercoribus}肥沃田地，\OriginalTerm{奥吉亚斯}{Augias}有更多的粪便可以施给田地。如果\OriginalTerm{皮库斯}{Picus}之子\OriginalTerm{法乌努斯}{Faunus}因疯狂而违犯法律和正义，那么他更适合被医治而非神化。如果\Person{法乌努斯}的女儿在贞洁上如此卓越，以至于不与男人交谈，那也许是因为粗鲁、自知丑陋或因父亲的疯狂而羞愧。比起这位\Quote{好女神}，\OriginalTerm{珀涅罗珀}{Penelope}更配得上神圣尊荣，尽管她身处众多最卑劣的求婚者之中，却以巧妙的手法保全了所遭攻击的纯洁。还有\OriginalTerm{桑克图斯}{Sanctus}，因其好客，\OriginalTerm{普洛提乌斯}{Plotius}王为他建了一座神庙；甚至\OriginalTerm{尤利西斯}{Ulysses}本可以把他赐予你们一位神，即最文雅的\OriginalTerm{阿尔基努斯}{Alcinous}。\par

% structure: p0020 source=h2 target=subsection reason=relative-heading-rank; base=h2

\subsection{第十章 罗马神话的一个可耻特征：它尊崇像拉伦提娜（Larentina）这样的臭名昭著人物。}

我急于提及更可憎的事例。你们的作家竟不羞于刊载有关\OriginalTerm{拉伦蒂娜}{Larentina}的事。她是一名雇来的妓女，或说是\OriginalTerm{罗慕路斯}{Romulus}的乳母，因此被称为\OriginalTerm{母狼}{Lupa}——因为她是个妓女；或说是如今已去世（即已神化）的\OriginalTerm{赫拉克勒斯}{Hercules}的情妇。他们讲述说，赫拉克勒斯的神庙看守人独自在庙中掷骰子；为了在无人陪伴时为自己找个对手，他开始一手代表赫拉克勒斯、另一手代表自己与自己对弈。（条件是：）若他从赫拉克勒斯处赢了赌注，便用这些钱置办一顿晚餐并雇一名妓女；但若赫拉克勒斯赢了（我指的是他的另一只手），则他应为赫拉克勒斯提供同样的东西。赫拉克勒斯的手赢了。这一成就本可被列入他的十二项功绩！看守人为这位英雄购买晚餐，并雇来拉伦蒂娜充当妓女。那曾焚化连赫拉克勒斯身体的火焰享用了晚餐，祭坛吞噬了一切。拉伦蒂娜独自睡在庙中；而她——一个出自妓院的女子——竟夸口在梦中屈从于赫拉克勒斯的快感；她或许在睡眠中确实经历了此事，因为这事曾掠过她的脑海。清晨，她早早离开神庙时，被一个年轻人搭讪——可谓\Quote{第三个赫拉克勒斯}。他邀请她去他家。她答应了，记起赫拉克勒斯曾告诉她这将对她有利。然后，他当然获准与她结为合法婚姻（因为若有人与神的姘妇交合而不受罚是绝不允许的）；丈夫立她为继承人。不久，在她临终前，她将经由赫拉克勒斯获得的那笔相当丰厚的遗产遗赠给罗马人民。之后，她也为女儿们寻求神化——她们本应被神圣的拉伦蒂娜立为继承人才是。罗马诸神因她的尊位而增添了一位。因为她在赫拉克勒斯的所有妻子中独受宠爱，只因她独自富有；她甚至比\OriginalTerm{凯雷斯}{Ceres}——那位曾为冥王（死者之王）提供欢愉的女神——更为幸运。在你们中间如此多的范例和\Emph{卓越}的名号之后，谁不曾被宣布为神呢？事实上，有谁曾对\OriginalTerm{安提诺乌斯}{Antinous}的神性质疑过？甚至\OriginalTerm{伽尼墨得}{Ganymede}比之于他，在爱他的（至高神）眼中，是否更为感恩和珍贵？照你们的说法，天堂向死者敞开。你们从冥府修筑了一条通往星辰的道路。妓女们从四面八方攀登上去，因此你们切勿以为自己正在赐予你们的国王莫大的荣耀。\par

% structure: p0022 source=h2 target=subsection reason=relative-heading-rank; base=h2

\subsection{第十一章　罗马人为出生——甚至在出生之前——直至死亡提供神祇。此体系中多有不宜。}

你们不仅满足于断言那些曾经为你们所知、你们曾听闻、触摸、其肖像被描绘、事迹被传述、记忆被保存于你们中间的人的神性；人们还坚持用一重天上的生命来祝圣我不知道什么无形的、无生命的影子，以及事物的\Emph{纯粹}名称——将人的整个存在从母腹中受孕起就分配给各自独立的能力：于是有一位\OriginalTerm{孔塞维乌斯}{Consevius}神，掌管交合生殖；一位\OriginalTerm{弗卢维奥娜}{Fluviona}，保护（子宫中）胎儿的（生长）；其后是\OriginalTerm{维图姆努斯}{Vitumnus}和\OriginalTerm{森提努斯}{Sentinus}，婴孩通过他们开始拥有生命和最初的感知；然后是\OriginalTerm{迪埃斯皮特}{Diespiter}，他负责使孩子完成出生。但当妇女开始分娩时，\OriginalTerm{坎德利费拉}{Candelifera}也\Emph{前来}相助，因为生育需要烛光；还有其他一些女神，她们的名字来自在产程中所扮演的角色。还有两位\OriginalTerm{卡尔门塔}{Carmentas}，按一般看法：一位名叫\OriginalTerm{波斯维尔塔}{Postverta}，其职责是协助足位胎儿的分娩；另一位叫\OriginalTerm{普罗萨}{Prosa}，则为正产婴儿执行同样的任务。\OriginalTerm{法里努斯}{Farinus}因（启发）最初发声而得名；另有人认为\OriginalTerm{洛库提乌斯}{Locutius}与言语的恩赐有关。\OriginalTerm{库尼娜}{Cunina}在场保护孩童的酣睡，为之提供安息之息。为了扶起（跌倒的孩童）有\OriginalTerm{莱瓦娜}{Levana}，还有与她在一起的\OriginalTerm{鲁米娜}{Rumina}。令人惊讶的疏忽是：没有指派任何神祇负责清除孩童的污秽。之后，主管他们最初的糊状食物和早期饮品的有\OriginalTerm{波蒂娜}{Potina}和\OriginalTerm{埃杜拉}{Edula}；教孩子直立是\OriginalTerm{斯塔蒂娜}{Statina}的工作，而\OriginalTerm{阿德奥娜}{Adeona}帮助他来到\Emph{亲爱的妈妈}跟前，\OriginalTerm{阿贝奥娜}{Abeona}则帮助他再次蹒跚走开；然后是\OriginalTerm{多米杜卡}{Domiduca}，（将新娘带回家）；还有\OriginalTerm{门斯}{Mens}女神，影响心智向善或向恶。他们还有\OriginalTerm{沃卢姆努斯}{Volumnus}和\OriginalTerm{沃莱塔}{Voleta}，掌管意志；\OriginalTerm{帕文蒂娜}{Paventina}，（恐惧之女神）；\OriginalTerm{维尼利亚}{Venilia}，希望之神；\OriginalTerm{沃卢皮亚}{Volupia}，快乐之神；\OriginalTerm{普雷斯提提亚}{Præstitia}，美丽之神。接着，他们为\OriginalTerm{佩拉格诺尔}{Peragenor}命名，因其教导人完成工作；为\OriginalTerm{孔苏斯}{Consus}命名，因其向人提供忠告。\OriginalTerm{尤文塔}{Juventa}是穿上成年长袍时的向导，而\Quote{长胡须的命运女神}则是在人达到成年时。若我必须提及他们婚配的职责，有\OriginalTerm{阿费伦达}{Afferenda}，其指定的职责是看管聘礼的献纳；但你们真该惭愧！你们有\OriginalTerm{穆图努斯}{Mutunus}和\OriginalTerm{图图努斯}{Tutunus}以及\OriginalTerm{佩尔通达}{Pertunda}和\OriginalTerm{苏比古斯}{Subigus}以及\OriginalTerm{普雷玛}{Prema}女神，同样还有\OriginalTerm{佩尔菲卡}{Perfica}。啊，放过你们自己吧，你们这些无耻的神祇！无人出席婚姻生活的隐秘挣扎。那极少数的、有那般意愿的人，离开并在极度的羞耻中面色绯红，就在他们的欢乐之中。\par

% structure: p0024 source=h2 target=subsection reason=relative-heading-rank; base=h2

\subsection{第十二章　最初的神祇是人——带有一些极为可疑的特征。\OriginalTerm{萨图尔努斯}{Saturn}或时间之神是人。关于他的观点相互矛盾。}

现在，我还有必要继续数说你们的神明吗？——因为我想详述你们所接纳的这类神明的特性。我实在不确定是该嘲笑你们的荒谬，还是该斥责你们的盲目。因为我要提出多少神明，而且是什么样的神明呢？是较大的，还是较小的？是古老的，还是新近的？是男性的，还是女性的？是未婚的，还是已婚的？是聪明的，还是拙笨的？是乡下的，还是城里的？是本国的，还是外来的？事实上，有如此众多的家族、如此众多的民族，它们需要一份神明名录，以至于根本无法考察、区分或描述它们。但主题越发散，我们就越必须加以限制。因此，在这次回顾中，我们只着眼于一个目标——证明所有这些神明曾经都是人（其实并非要教导你们这个事实，因为你们的行为表明你们已经忘记了它）——让我们从最自然的考察方法中采用一个简要的总结，甚至考虑他们种族的起源。因为起源决定了其后的一切。按我的看法，你们神明的这种起源始于萨图尔努斯。当瓦罗提到朱庇特、朱诺和密涅瓦是最古老的神时，我们本不该忽略这一点：任何父亲都比他的儿子更古老，因此萨图尔努斯必定先于朱庇特，就像卡埃卢斯先于萨图尔努斯一样，因为萨图尔努斯出自卡埃卢斯和特拉。不过，我且略过卡埃卢斯和特拉的起源。他们以某种不可思议的方式过着独身生活，没有子女。当然，他们需要很长时间才能健壮地长到那样的身高。不久，当卡埃卢斯的嗓音开始变化，特拉的胸脯变得结实，他们便彼此结为婚姻。我想，要么是天降下到他的配偶那里，要么是地上腾去迎接她的夫君。无论怎样，特拉受了天的种子，满期后奇妙地生下了萨图尔努斯。他像父母中的哪一位呢？那么，即使在那之后开始有了父母，可以肯定的是，在萨图尔努斯之前他们没有孩子，之后只有一个女儿——俄普斯；从此他们不再生育。事实上，萨图尔努斯趁卡埃卢斯睡觉时阉割了他。我们读到这个名称\Quote{卡埃卢斯}是阳性。况且，他若不是男性，又怎能作父亲？但阉割是用什么工具进行的？他有一把镰刀。怎么，那么早就有镰刀了？因为武尔坎当时还不是铁匠。然而，守寡的特拉虽然还很年轻，却并不急于再婚。确实，没有第二个卡埃卢斯给她。除了俄刻阿诺斯，谁还向她投怀送抱？但他带着咸味，而特拉习惯的是淡水。因此，萨图尔努斯是卡埃卢斯和特拉唯一的儿子。他长到青春期，娶了自己的妹妹。当时还没有法律禁止乱伦或惩罚弑亲。然后，当儿子们出生后，他会吞吃他们；与其让狼吃掉他们（如果他遗弃他们，他们就会成为狼的食物），不如他自己吞掉。无疑，他害怕某个儿子会学到父亲那把镰刀的教训。后来朱庇特出生时，他被挪走了：父亲吞下了一块石头代替儿子，这是假装的。这个诡计暂时保护了他的安全；但最终，那个他未曾吞吃、在秘密中长大的儿子袭击了他，剥夺了他的王国。这就是神明们的始祖，是天地为你们所生，有诗人作助产士。现在，一些想象力丰富的人认为，通过这个关于萨图尔努斯的寓言故事，对\Emph{时间}做了一个生理学的描绘：（他们认为）因为万物都被时间摧毁，所以卡埃卢斯和特拉本身是父母却没有任何自己的父母，并且使用了那把（致命的）镰刀，而且（萨图尔努斯）吞吃自己的后代，因为实际上他吸收了一切从他自己产生的东西。他们还援引他的名字为证；因为他们说，希腊语中他被称为\OriginalTerm{克罗诺斯}{\GreekText{Κρόνος}}，意思与\OriginalTerm{时间}{\GreekText{χρόνος}}相同。他的拉丁语名称他们也从\Emph{种子}衍生而来；因为他们认为他实际上是万物的创造者——正是通过他种子从天上降落到地上。他们将他与\Emph{俄普斯}联合，因为种子产生了实际生命的丰盛\Emph{财富}，并且因为种子通过\Emph{劳动}发育。现在，我希望你们解释这个比喻的说法。它要么是萨图尔努斯，要么是时间。如果是时间，它怎能是萨图尔努斯？如果是萨图尔努斯，它怎能是时间？因为你们不可能将这两个有形的实体视为同一个人。然而，有什么阻止你们按照时间的适当性质来崇拜时间呢？为什么不让一个人，甚至一个神话中的人，成为你们崇拜的对象，但各自按照其本性，而非以时间的性质？你们这种智力的造作有何意义，如果不是用似是而非的理由的表象来粉饰最污秽的事？另一方面，你们既不打算将萨图尔努斯当作时间，因为你们说他是一个人；也不打算在将他描绘成时间时，因此而认为他曾是人。无疑，在远古的记载中，你们的神明萨图尔努斯被明确描述为以人形生活在地上。任何从未存在过的东西显然都可以被描绘成无形的；哪里有真实，哪里就没有这种虚构的余地。因此，既然有明确的证据表明萨图尔努斯曾经存在，你们改变他的特性是徒劳的。你们不能否认他曾是人，他就不任凭你们随意处置，也不能断言他要么是神圣的，要么是时间。在你们每一页文献中，萨图尔努斯的起源都很显眼。我们在卡西乌斯·塞维鲁、科尔内利乌斯们——内波斯和塔西佗那里读到，在希腊人中也在狄奥多罗斯以及所有其他古代编年史编纂者那里读到。关于他，没有比在意大利本身更可靠的记录了。因为，经过许多国家，并受到雅典的款待之后，他定居在意大利，或称为埃诺特里亚，受到雅努斯——或者像萨利祭司称呼的那样，雅纳斯——的热情欢迎。他定居的山丘被称为萨图尔尼乌斯，他所建立的城市至今仍称为萨图尔尼亚；简言之，整个意大利曾一度有同一名称。这就是来自现在世界主宰的那个国家的见证：关于萨图尔努斯的起源，无论有何疑问，他的行动清楚地告诉我们他曾是人。因此，既然萨图尔努斯是人，他无疑来自人的血统；而且，因为他是人，他当然不是来自卡埃卢斯和特拉。然而，有些人觉得很容易称他为那些有时似乎万物都从之而来的神的儿子，正如父母不详的人一样。因为谁不是在敬畏中称天和地为自己的父亲和母亲呢？或者，按照人的习惯，当提到任何不知名或突然出现的人时，他们说\Quote{他们是从天上来的吗？}因此，当一个陌生人到处突然出现时，人们习惯称他为天生之人——就像我们也常称所有家世不明的人为地生之人一样。我不提这样一个事实：上古时代如此，人们的眼睛和思想都习惯于粗陋，以至于每个新来者的出现都像神一样令他们激动；对于一位君王，尤其是原始君王，更是如此。我将在萨图尔努斯这个例子上多停留一会儿，因为通过充分讨论他的原始历史，我将预先为所有其他情况提供一个简要的答案；而且我不想省略你们神圣文献中更有说服力的证据，这些证据的可信度应该与其古老程度成正比。现在，比所有文献更古老的是西比拉；我指的是那位真正的真理女先知，你们借用她的名字来称呼你们魔神的祭司。她用六音步诗行阐述了萨图尔努斯的世系及其事迹，话语如下：\Quote{在洪水淹没前代之后的第十代，统治着萨图尔努斯、提坦和伊阿珀托斯，特拉和卡埃卢斯之子中最勇敢的。}因此，无论对你们较古老的作者和文献有多少信任，而且对那些属于那个时代的最纯朴的作者和文献更加信任，就足以确定萨图尔努斯及其家人是人。于是，我们拥有一个简明的原则，它对于所有其他情况来说相当于一条关于他们起源的惯例规则，以避免我们在个别情况中出错。后裔的特殊性格由种族的原始奠基者显示——凡人来自凡人，地上的人来自地上的人；一步一步依次而来——婚姻、怀孕、出生——国家、定居点、王国，一切都提供最清晰的证明。因此，那些不能否认人生的人，也必须承认人死；那些承认人有死性的人，不应认为他们是神。\par

% structure: p0026 source=h2 target=subsection reason=relative-heading-rank; base=h2

\subsection{第13章——诸神最初是人：谁有权使他们成为神？\Person{尤比特}不仅是一个人，而且是不道德的。}

像\Person{瓦罗}及其同梦者这样的人，将那些他们无法断言在原始状态中不是人的人纳入神明的行列；（他们这样做）是断言那些人死后成为了神。那么，我在此持定立场。如果你们的神明是被选任到这种尊荣和神性，就像你们招募元老院成员一样，那么你们不得不承认（以你们的智慧），必定有一位至高无上的主宰，有权挑选，如同一位\Person{凯撒}；没有人能将他没有绝对控制权的东西赐予他人。此外，如果他们死后能自己成神，请告诉我，为何他们起初选择处于一个较低的状态？或者，如果没有谁使他们成神，他们又如何能被说成是被如此造成，既然他们只能由别人造成？因此，你们没有理由否认存在某位批发神性的主宰。让我们接着审视将凡人送往天堂的理由。我想你们会提出一对理由。那么，无论谁（授予神圣荣誉），他行使职能，要么是为了获得一些支持、防御，甚至是他自己尊严的装饰；要么是基于有功劳者的迫切要求，以奖赏所有应得之人。我们不允推测其他原因。没有人赠予他人礼物时不考虑自己的利益或对方的利益。然而，这种行为不配给天主，因为祂的能力如此巨大，可以直接造神；而祂借口人需要甚至死者的帮助和支持，把人抬举到如此高的地位，这是一种奇怪的念头，因为祂从起初就能为自己创造不朽的存在。已比较人性与神性之事的人无需在这些点上更多论证。然而，后一种意见应当被讨论，即天主因功劳的考量而授予神圣荣誉。那么，如果授予是基于这样的理由，如果天堂因为古人的功德而向他们敞开，我们必须思考，此后没有谁配得这样的荣誉；除非如今再也没有这样的地方可以让人达到。让我们假设古时人们可能因巨大的功德而配得天堂。然后让我们考虑是否真的有这样的功德。让主张存在功德的人宣告他自己的功德观。既然人类在时间初期所做的事是他们被神化的有效主张，你们一贯地将荣誉授予犯下乱伦罪过的兄妹——\Person{俄普斯}和\Person{萨图尔努斯}。你们的\Person{尤比特}，在婴儿时被偷走，不配受到人类所得到的家和滋养；正如他作为坏孩子所应得的，他不得不生活在\Place{克里特岛}。后来，当长大成人，他废黜了自己的父亲，无论他父亲的为人如何，他的统治最为繁荣，因为他是黄金时代的国王。在他统治下，远离劳苦和匮乏，和平保持着欢乐而温和的统治；在他统治之下——\par

\Quote{没有农夫征服田地；}\newline{}\Quote{没有牧人使田地臣服于他们的权下；}\par

没有人的催促，大地就自动长出所有庄稼。但他恨一个犯过乱伦、曾阉割自己祖父的父亲。然而，看哪，他自己娶了自己的妹妹；所以我想那句古语是为他而说的：\OriginalTerm{有其父必有其子}{\GreekText{Τοῦ} \GreekText{πατρὸς} \GreekText{τὸ} \GreekText{παιδίον}}——\Quote{父亲的亲孩子}。在父亲的虔诚和儿子的虔诚之间\Quote{没有差别}。如果即使在那早期法律是公正的，\Person{尤比特}应当被\Quote{装进两个麻袋缝起来}。在这桩通过乱伦满足的淫欲得到证实之后，他为何还要犹豫放任自己沉溺于通奸和放荡这些较轻的过犯？自从诗歌如此戏弄他的品性，就像通常张贴逃亡奴隶的告示一样，我们习惯于在与路人闲谈时无拘无束地议论他的诡计；有时将他描绘成正是他淫乱代价的那笔钱的形式——例如当他（扮演）一头公牛，或者不如说付了一头公牛的价，并将（金子）倾泻进少女的闺房，或者不如说用贿赂硬闯进去；有时则描绘成他所扮演角色的形象——例如掠夺（美少年）的鹰，以及歌唱（迷人歌曲）的天鹅。那么，如此寓言岂非由最令人作呕的阴谋和最坏的丑闻构成？人们的风俗和性情岂非不会因这些榜样而变得放荡？魔鬼——长久从事其使命的恶天使的后裔——如何努力使人离开信德转向不信和如此寓言，我们在此不应过多谈论。的确，你们神明的整体以他们的君王、王公和导师为榜样，并非同一性质；是以某种其他方式，他们的权威要求品性相似。但他们中最坏的他，岂非（本应但）不是他们中最好的？凭着对他独有的称号，你们习惯称\Person{尤比特}为\Quote{至善}，而在\Person{维吉尔}笔下，他是\Quote{公正的尤比特}。因此，所有的人\Emph{像}他——对自己的亲属乱伦，对陌生人淫乱，不敬神，不义！如今，神话故事未留下明显恶名的他，不配成为神。\par

% structure: p0030 source=h2 target=subsection reason=relative-heading-rank; base=h2

\subsection{第14章——被承认被提升到神性境况的神明，他们有何突出的权利获得这样的荣誉？\Person{赫丘利}是一个较低等的人物。}

然而，既然他们坚持认为，凡由人类状态被接纳至神化尊荣者应与其他人分开，且狄奥尼修·斯多噶所划定的本有之神与人造之神应当加以区分，那么我也愿就最后一类略赘数言。我将以\Person{赫丘利}本人为例，引出答问之要旨：他是否堪配天堂与神性尊荣？因为照人类所愿，这些尊荣是按功德赐予他的。若因其无畏消灭野兽的勇毅而予之，这其中又何等值得铭记？难道被判刑而投入竞技场的罪犯，即使被交付于卑劣角斗场的角逐，不也常一次杀死数头野兽，且更为热切么？若因其周游世界，阔绰者于愉快闲暇之中，或哲学家于奴隶般的贫困之中，岂非同样屡次完成此等旅程？难道忘了犬儒\Person{阿斯克勒庇阿得斯}仅凭一头可怜母牛，骑其背上，有时吮其乳汁，便亲身走遍全球？即便\Person{赫丘利}曾巡游阴间，谁不知往哈得斯之路对众人皆开？若因其屡次屠杀与众多战役而奉为神明，那么著名的\Person{庞培}，那位征服了连\Place{奥斯提亚}亦遭其劫掠的海盗之人，岂非赢得更多胜利？再问我，多少千人曾被围于\Place{迦太基}卫城一角，为\Person{西庇阿}所杀？因此，\Person{西庇阿}比\Person{赫丘利}更有资格被视为神化的候选者。你更须留意，将\Person{赫丘利}与妾妇及妻子的淫乱、\Person{翁法勒}的束带、以及因失去俊美少年而可耻地抛弃\OriginalTerm{阿尔戈英雄}{Argonauts}等行径，也加于其功德之上。在这可耻的标记上，再为其荣光增添癫狂发作，敬拜那杀死其子女与妻子的箭矢。此人，在悔恨弑亲之苦中自以为堪配火葬柴堆，最终更应披上其妻因他好色而送来的毒袍，受那应得的可耻之死。然而，你们却以同样轻易的态度，将他从火堆升上天际，如同（你们同样）荣耀了另一位亦为神火暴力所灭的英雄。此人发明了少许实验，据称能借医术使死者复生。他是\Person{阿波罗}之子，半属人类，虽为\Person{朱庇特}之孙、\Person{萨图尔努斯}之曾孙（实则血统可疑，因其出身不明，如\Person{阿尔戈}的\Person{苏格拉底}所述；他亦遭遗弃，受比朱庇特更糟的监护，竟由母狗乳头喂养）；无人能否认他遭雷击而死乃其应得之终。在此事中，你们最卓越的\Person{朱庇特}再次犯错——对孙儿不敬，嫉其技艺。\Person{品达}其实并未隐瞒其真实功过；据他所说，此人因贪财好利而受罚，受此驱使，他滥用可出售的医术，以反常方式令生者死去，而非令死者复生。据说其母亦遭同一击而死，她将这头危险野兽赐予世界，理应借同一阶梯逃往天堂。雅典人仍将不乏如何向此类神明献祭之道，因他们在亡者中向\Person{埃斯库拉庇俄斯}及其母致以神性尊荣。仿佛他们手边没有自己的\Person{忒修斯}可供崇拜，他如此堪配神之殊荣！为何不？他岂非于异乡遗弃救命恩人，其冷漠乃至无情，正如其成为父亲之死因一般？\par

% structure: p0032 source=h2 target=subsection reason=relative-heading-rank; base=h2

\subsection{第十五章  星座与守护神均为极无价值之神祇。罗马对神祇的垄断令人不满。其他民族同样需要神灵。}

逐一审视所有那些被你们葬于星座之中、并胆敢奉为神灵者，实为冗长。我想你们的卡斯托尔、珀尔修斯、厄里戈娜，与\Person{朱庇特}自己的大男孩有同样资格获得天界尊荣。但何足为奇？你们甚至将狗、蝎子、蟹都迁至天上。关于你们在神谕中所拜者，我暂不置评。此崇拜之存在，由宣示神谕者所证实。为何？你们甚至会有神祇成为悲伤的旁观者，如\Person{维杜斯}，他使灵魂成为\Emph{寡妇}，将灵魂与身体分离，而你们判他不得被围入城垣之内；还有\Person{卡库鲁斯}，剥夺眼睛的感知；\Person{奥尔巴娜}，使种子丧失生机；此外还有死亡女神本人。匆匆略过其余，你们将地方或城市之场所奉为神祇，如\Person{雅努斯}神父（此外还有弓女神\Person{雅娜}），以及七丘之\Person{塞普提蒙提乌斯}。\par

人们在同一个地方设有祭台或庙宇，向相同的\OriginalTerm{吉尼乌斯}{Genii}献祭；但当他们寄居异乡或租住房屋时，也向其他\OriginalTerm{吉尼乌斯}{Genii}献祭。我且不提\OriginalTerm{阿森苏斯}{Ascensus}，其名源自攀登的癖好，以及\OriginalTerm{克里维科拉}{Clivicola}，其名源自斜坡的居所；我也不提那些被称为\OriginalTerm{福尔库卢斯}{Forculus}（来自门扇）、\OriginalTerm{卡尔德亚}{Cardea}（来自门铰）、\OriginalTerm{利门提努斯}{Limentinus}（门槛之神）的神祇，以及你的邻居们敬拜的其他街门守护神。这并不奇怪，因为人们在妓院、厨房甚至监狱中都有各自的神祇。因此，天上挤满了无数属于它自己的神祇，既有这些，也有属于罗马人的其他神祇，他们将人一生的各种职能分配给他们，以至于并不缺乏其他神祇。诚然，我们所列举的神祇被认为是罗马独有的，在异域不易被承认；然而，在整个人类种族和每个民族中，那些人意欲其神祇掌管的一切职能和情境又是如何出现的呢？在这些情境中，他们的保证不仅没有得到官方承认，甚至根本没有得到任何承认。\par

% structure: p0035 source=h2 target=subsection reason=relative-heading-rank; base=h2

\subsection{第十六章　实用技艺的发明者不配被神化。他们必先承认造物主。技艺因时而变，有些已过时。}

然而，有些人发现了果实和生活中各种必需品，因而配得被神化。现在请问，当你们称这些人为\Quote{发现者}时，你们岂不是承认他们所发现的东西早已存在了吗？那么，你们为何不更愿意尊崇那位真正赐予恩赐的创造主，反而将创造主转变为仅仅是发现者呢？先前，那作出发现的人，即发明者本人，无疑向创造主表达了感恩；无疑，他也感到他就是天主，宗教敬拜本当属于他，因为他是恩赐的创造者，发现者以及所发现之物都是由他创造的。阿非利加的绿无花果，在加图将其引入元老院时，罗马无人知晓，他是为了表明那敌人的行省多么近，他不断敦促要征服该行省。樱桃首次在意大利普及，是由格奈乌斯·庞培从本都引入的。我或许曾认为，最早将苹果引入罗马的人，配得公开神化的荣誉。然而，这作为制造神的理由，与发明实用技艺一样愚蠢。然而，若将我们时代的能工巧匠与这些人相比，那么后来的世代岂不是比前者更适合被神化？请告诉我，所有的现存发明不都超越了古代吗？而日常经验不断增添新的库存。因此，你们因他们的技艺而视为神明的人，你们实际上是在用你们的技艺伤害他们，并用无法超越的竞争技艺挑战他们的神性。\par

% structure: p0037 source=h2 target=subsection reason=relative-heading-rank; base=h2

\subsection{第十七章　结论：罗马人的帝国权力并非来自其神祇。唯有伟大的天主赐予王国；他是基督徒的天主。}

总而言之，我并未否认古人愿意崇敬、后代相信为神祇的一切都是你们宗教的守护者；但仍有待我们考虑的是，罗马迷信中一个非常大的假设必须在我们反对你们（异教徒）时加以讨论，即：罗马人之所以成为全世界的君主与主宰，是因为他们通过宗教职事赢得了如此统治权，以至于他们几乎在权能上超越了他们的神祇。难怪\OriginalTerm{斯特尔库卢斯}{Sterculus}、\OriginalTerm{穆图努斯}{Mutunus}和\OriginalTerm{拉伦提娜}{Larentina}各自将这个帝国推向了巅峰！罗马人民唯独由他们的神祇指定获得如此统治权。因为我无法想象任何外邦神祇会宁愿为一个陌生民族所做之事多于为自己的人民，从而成为其出生、成长、显赫和埋葬的祖国的逃兵、疏忽者，甚至是背叛者。因此，就连\OriginalTerm{朱庇特}{Jupiter}也不能容忍他自己的\Place{克里特}被罗马的权杖征服，忘记那\Place{伊达}山洞、\OriginalTerm{科律班忒斯}{Corybantes}的铜钹，以及在那片亲爱之地哺育他的山羊的宜人气味。他难道不会使他的坟墓超越整个\Place{卡皮托利}，以便那覆盖朱庇特骨灰的土地最广泛地统治吗？同样，\OriginalTerm{朱诺}{Juno}会愿意那因她的爱而甚至忽略了\Place{萨摩斯}的布匿城被毁灭，而且是被\Person{埃涅阿斯}子孙的火焰所毁灭吗？虽然我深知：\par

\Quote{此处有她的武器，\newline{}此处有她的战车，她早已渴望并抚育着\newline{}这王国应归于万民，\newline{}如果命运允许的话。}\par

\Quote{这里有她的武器，这里有她的战车，\newline{}此处女神般，为将一日\newline{}普世治权之座固定，\newline{}倘若命运能被扭转为许可，\newline{}甚至那时她的计划，她的挂念已倾注。}\par

\Emph{然而}不幸的（众神之后）竟无力反抗命运！但罗马人并没有给予命运女神同等的尊荣，尽管他们将迦太基交在她们手中，却给了拉伦提娜更多。然而，你们那些神祇确实没有赐予帝国的能力。因为当朱庇特在克里特为王，萨图尔努斯在意大利，伊西斯在埃及时，他们是以人的身份统治，并且有许多人协助他们。因此，那服役者（阿波罗）也使他人成为主人，阿德墨托斯的奴隶（阿波罗）以帝国扩张了罗马公民的权势，尽管他用模棱两可的神谕欺骗了自己慷慨的信徒克洛伊索斯，致使他亡国。身为神祇，他为何不敢大胆预言他将失去王国的真相？的确，那些被赋予执掌帝国权柄之神，理应能时刻守护自己的城市。若他们有能力将帝国赐予罗马人，为何密涅瓦没有守护雅典免遭薛西斯之害？为何阿波罗没有将德尔斐从皮洛士手中救出？他们既丧失了自己的城市，却保存了罗马城，因为罗马的虔诚（据说）配得保护！但事实难道不是：这过度的虔诚是在帝国通过权力增长获得荣耀之后才被杜撰出来的吗？无疑，神圣礼仪是由努马引入的，但那时你们的行事尚未被偶像和神庙的宗教所败坏。虔诚是单纯的，崇拜是谦卑的；祭台简陋，器皿朴素，乳香稀少，神祇本身无处可寻。因此，人并非先伟大而后虔诚，也非因虔诚而伟大。然而，罗马人如何可能看似因过度的虔诚和对神祇的极大尊敬而获得帝国，既然帝国反而是在神祇被轻视之后才扩大的？如果我没有弄错，每一个王国或帝国都是通过战争获得和扩大的，而征服者同样伤害了它们的神祇。因为同样的毁灭临到城墙和神庙；平民和祭司遭遇同样的屠杀；世俗物品和神圣物品遭受同样的掠夺。罗马人拥有多少战利品，就犯下多少渎神之罪；他们对神祇的胜利与对民族的胜利一样多。被俘的偶像至今仍在他们当中；当然，若它们能看到征服者，绝不会爱他们。然而，既然它们没有知觉，被伤害也安然无恙；既然被伤害安然无恙，崇拜它们便毫无益处。因此，一个通过一场场胜利强大起来的民族，不可能看似因宗教的功德而发展——无论是他们因增强权力而损害了宗教，还是因损害宗教而增强了权力。所有民族都曾在其特定时期拥有帝国，如亚述人、米底人、波斯人、埃及人；帝国至今仍为一些民族所拥有，然而，那些失去权力者，即使在不蒙神祇眷顾之后，也未尝不重视宗教仪式和敬拜神祇，直到最后几乎普世的统治权落于罗马人之手。这是时代的命运，不断以革命动摇王国。询问谁规定了这些时代的变化。是那同一位伟大存在，他分配王国，并已将最高统治权置于罗马人手中，如同多国的贡品被征收后集于一个大箱。他对此有何决定，那些最亲近他的人知道。\par

% structure: p0042 source=h2 target=subsection reason=relative-heading-rank; base=h2

\subsection{附录。一篇关于外邦人可憎神祇的残篇。}

\EditorialNote{欧勒指出此残篇为伪作，认为它出自一个对特土良风格与教导仅略知一二者之手。我未在杜平或鲁思的著作中找到提及。本译文由塞尔沃提供。}\par

如此巨大的盲目已降临罗马民族，以致他们称自己的仇敌为\Quote{主}，宣扬恩赐的窃贼即是赐予者本人，并向他献上感谢。于是，他们用人的名字称呼那些神祇，而非神祇自己的名字，因为他们不知神祇的名字。他们明白那些是魔鬼：但阅读古代诸王的史书后，尽管看到他们的品性是可朽的，仍以神的名号尊崇他们。\par

至于他们称之为\Person{朱庇特}、并认为是最高之神的那位，他出生时，从世界创立到他已有大约三千年。他生于\Place{希腊}，父母是\Person{萨图尔努斯}和\Person{奥普斯}；因怕被他父亲杀死（或者，如果可以这么说，怕被重新生出），在他母亲建议下被带到\Place{克里特}，在\Place{伊达山}的洞穴中养大；靠\Place{克里特}人——凡胎所生！——敲击兵器而躲过父亲的搜寻；吸吮母羊的奶；剥下羊皮；披上羊皮；他亲手杀死了自己的乳母，并用她的皮！（据作者\Person{荷马}所述，如果可信的话，他还在上面缝了三个金穗，每个价值一百头牛。）这位\Person{朱庇特}成年后与父亲交战多年；打败了他；对家宅进行弑父式的袭击；奸污了自己的童贞姐妹；选了其中一位结婚；用武力驱逐了父亲。此外，那场行动的其他场景也有记载。他通过他人的妻子或被奸污的童贞女子生了儿子；玷污了自由民男孩；以暴君和君主的统治不法地压迫各民族。他们错误地以为是\Emph{原初}神的父亲，竟不知道这个儿子藏在\Place{克里特}；而他们相信是\Emph{更强大}神的儿子，也不知道被他放逐的父亲潜藏在\Place{意大利}。若他在天上，怎能看见\Place{意大利}的事呢？因为\Place{意大利}之地并不在\Quote{角落}。然而，若他是\Emph{神}，本不应有任何事瞒过他。但意大利人称其为\Person{萨图尔努斯}的那位确实藏在那里，这从以下事实便可清楚看出：因他藏匿，\Place{赫斯佩里亚}的语言至今仍被称为拉丁语，正如他们的作者\Person{维吉尔}所述。于是，据说（\Person{朱庇特}）生于地上，而（他父亲\Person{萨图尔努斯}）害怕被他逐出王国，想杀死他这个竞争者，却不知道他已被秘密带走并藏了起来；后来，子神追赶父亲，不朽者想杀死不朽者（这可信吗？），却因一片海隔而未能得逞，且不知（猎物的）逃跑；当两个神之间在地上发生这一切时，天空被抛弃了。无人降下雨水，无人打雷，无人管理这整个世界的重负。因为他们甚至不能说他们的行动和战争发生在天上；这一切都发生在\Place{希腊}的\Place{奥林匹斯山}上。然而，天空不叫奥林匹斯，天空就是天空。\par

这些，就是他们的行为，我们将首先处理——出生、藏匿、无知、弑父、通奸、猥亵——这些不是神所行的，而是最不洁、最凶残的人所行的；这些人若活在今天，本应受所有法律的指控——这些法律远比他们自己的行为更公正、更严厉。\Quote{他用武力驱逐了父亲。}法尔基迪安法和森普罗尼安法会将弑亲者与野兽一起缝入袋中。\Quote{他奸污了姐妹。}帕皮尼安法会以所有刑罚处罚这种暴行，将肢体逐一惩治。\Quote{他侵犯了别人的婚姻。}尤利安法会以死刑惩治通奸的破坏者。\Quote{他玷污了自由民男孩。}科尔内利安法会以新的严刑处罚违犯性结合之罪，因其对新的结合犯下亵渎之罪。这个人被证明也不具有神性，因为他是一个人；他父亲的逃走瞒过了他。对这样一个凡人，这样一个邪恶的国王，如此淫荡而残忍，人竟将神的尊荣给了他。如今，当然，如果他\Emph{在地上}出生，通过生命各个时期的发展而成长，并在其中行了所有这些恶事，却不再在地上，那人们除了以为他死了还能想什么？或者，愚蠢的错误以为在他年老时生出了翅膀，好让他飞向天上？甚至\Emph{这}也可能在失去理智的人中间获得信任，如果他们真的相信（如他们所信的）他变成一只天鹅，生下\Person{卡斯托尔}；一只鹰，玷污\Person{伽尼墨得}；一头公牛，奸污\Person{欧罗巴}；黄金，奸污\Person{达那厄}；一匹马，生下\Person{皮里托俄斯}；一只山羊，从一只母羊生下\Person{埃吉帕}；一个\Person{萨提尔}，拥抱\Person{安提奥佩}。目睹这些通奸行为（罪人易犯），他们因此轻易相信，不端行为和一切污秽的制裁都是从他们虚构的神那里借来的。他们是否看到，他生涯中其余那些可以相信的行为（其实是真实的，且他们说他没有变形就行）是多么毫无改进？从\Person{塞墨勒}生了\Person{利伯尔}；从\Person{拉托娜}生了\Person{阿波罗}和\Person{狄安娜}；从\Person{玛亚}生了\Person{墨丘利}；从\Person{阿尔克墨涅}生了\Person{赫拉克勒斯}。至于其余那些他们自己承认的败坏，我不愿记录，以免已埋葬的丑恶再度回响在人耳边。但我已提到这些少数的后代；这些后代，他们在错误中也相信他们自己也是神——即生于一个乱伦的父亲；是通奸所生，是假冒的。而永生、永恒的天主，具有永恒的神性，预知未来，不可测量，他们却通过如此不可言说的罪行（将祂）消散（为无有）。\par

\SourceRecord{Translated by Peter Holmes.；Ante-Nicene Fathers；Vol. 3.；Edited by Alexander Roberts, James Donaldson, and A. Cleveland Coxe.；Buffalo, NY: Christian Literature Publishing Co.,；1885.；<http://www.newadvent.org/fathers/03062.htm>.}
